%% 3-7 反比例のグラフ（第12回：新しく書いた。組 V3-4）  % noruby
%% 2-15 の「くわしくは 3-7」を受ける：軸に近づくがくっつかない・直線でない・原点を通らない。  % noruby
%% 道具：①点をかく ②軸にくっつかないわけ（かけ算・面積の2通り）③x が負のがわ（双曲線）④a が負 ⑤変化の割合 ⑥pV＝一定  % noruby
\makeatletter
% 先をふさいだ注射器（横向き）  \bkthrhsyringe{(左下)}{空気の長さ mm}  つつの長さ50 mm・高さ10 mm  % noruby
\newcommand\bkthrhsyringe[2]{\begin{scope}[shift={#1}]
  \fill[bkGray!60] (-4,3.6) rectangle (0,6.4);
  \fill[bkRed!75] (-7,2.4) rectangle (-4,7.6);
  \fill[bkBlue!14] (0,0) rectangle (#2,10);
  \foreach \bkthrh@k in {1,...,7} {\fill[bkBlue!45] ({#2*(\bkthrh@k-0.5)/7},{5+2.6*sin(\bkthrh@k*137)}) circle (0.55);}
  \fill[bkInk!70, rounded corners=0.4mm] (#2,0.3) rectangle ({#2+3},9.7);
  \fill[bkGray!45] ({#2+3},4) rectangle ({#2+45},6);
  \fill[bkGray!70, rounded corners=0.5mm] ({#2+45},-2) rectangle ({#2+47.5},12);
  \draw[bkGray!90!black, line width=0.8pt] (0,0) -- (50,0) (0,10) -- (50,10) (0,0) -- (0,10);
  \draw[bkGray!90!black, line width=1.4pt] (50,-3) -- (50,13);
  \foreach \bkthrh@k in {1,...,9} {\draw[bkGray, line width=0.4pt] (\bkthrh@k*5,10) -- (\bkthrh@k*5,8.6);}
\end{scope}}
% 式を、まん中に1行で置く  \bkthrhline{…}  % noruby
\newcommand\bkthrhline[1]{\par\addvspace{3pt}{\centering #1\par}\addvspace{3pt}}
% 説明の段の見出し（丸数字＋太字の1行）  \bkthrhstep{番号}{見出し}  % noruby
\newcommand\bkthrhstep[2]{\par\addvspace{6pt}\Needspace{4\baselineskip}\noindent
  {\bkcircled[bkOrange]{#1}\enspace\sffamily\bfseries\color{bkOrange!75!black}#2}\par\nobreak\vspace{2pt}}
% 本文の中の小見出し（それだけで1行）  \bkthrhhead{見出し}  % noruby
\newcommand\bkthrhhead[1]{\par\addvspace{4pt}\Needspace{3\baselineskip}\noindent{\sffamily\bfseries #1}\par\nobreak\noindent}
\makeatother

\bkchapter{3}{7}{\R{反比例}{はん|ぴ|れい}のグラフ}{\bklevel{2}}{\R{反比例}{はん|ぴ|れい}のグラフ}
\bkmichishirube{2-15,3-2,3-4,3-5}{3-11,5-4,5-14}

\begin{goal}
\begin{itemize}[label={\color{\bkTT}●}, leftmargin=1.6\zw, itemsep=0pt, topsep=1pt]
\item \bklevel{2}\enspace $y=\dfrac{a}{x}$ の\R{表}{ひょう}を\R{作}{つく}って\R{点}{てん}をとり、なめらかな\R{曲線}{きょく|せん}（\R{双曲線}{そう|きょく|せん}）がかける。\R{軸}{じく}にくっつかないわけを\R{説明}{せつ|めい}できる。
\item \bklevel{2}\enspace $x$ や $a$ が\R{負}{ふ}のときもふくめてグラフの\R{形}{かたち}を\R{読}{よ}み、\R{変化}{へん|か}の\R{割合}{わり|あい}が\R{一定}{いっ|てい}でないことを\R{確}{たし}かめられる。
\item \bklevel{2}\enspace $xy$ が\R{一定}{いっ|てい}の\R{関係}{かん|けい}（$pV$ が\R{一定}{いっ|てい}、てこ、\R{速}{はや}さと\R{時間}{じ|かん}）を\R{見}{み}つけ、\R{式}{しき}を\R{作}{つく}って\R{使}{つか}える。
\end{itemize}
\end{goal}

\begin{hatena}
\noindent\begin{minipage}[c]{0.5\linewidth}
\R{注射器}{ちゅう|しゃ|き}に\R{空気}{くう|き}を\R{入}{い}れ、\R{先}{さき}をゴムのふたでふさいで、ピストンをおしこみます。おしこむほど、\R{手}{て}をおし\R{返}{かえ}す\R{力}{ちから}は\R{強}{つよ}くなります。

\R{空気}{くう|き}を\R{半分}{はん|ぶん}までおしちぢめると、おし\R{返}{かえ}す\R{強}{つよ}さは\R{何倍}{なん|ばい}になるかな？　もっとおしこんで、\R{空気}{くう|き}を0にすることはできるかな？　この\R{関係}{かん|けい}をグラフにかくと、どんな\R{形}{かたち}になるでしょう。
\end{minipage}\hfill
\begin{minipage}[c]{0.47\linewidth}\centering
\begin{tikzpicture}[x=0.62mm, y=0.62mm]
  \bkthrhsyringe{(0,22)}{40}
  \bkthrhsyringe{(0,0)}{20}
  \node[font=\sffamily\footnotesize, anchor=south west, inner sep=1pt] at (-7,35) {はじめ（\qty{20}{mL}）};
  \node[font=\sffamily\footnotesize, anchor=south west, inner sep=1pt] at (-7,13) {\R{半分}{はん|ぶん}（\qty{10}{mL}）};
  \draw[bk arrow, bkRed] (78,5) -- (70,5);
  \node[font=\sffamily\bfseries\Large, text=bkOrange!80!black] at (31,-6.5) {？\R{倍}{ばい}};
\end{tikzpicture}
\end{minipage}
\end{hatena}

2-15では、$xy$ が\R{一定}{いっ|てい}になる2つの\R{量}{りょう}を\R{反比例}{はん|ぴ|れい}とよび、そのグラフが「\R{軸}{じく}に\R{近}{ちか}づくけれど、くっつかない\R{曲線}{きょく|せん}」になることを\R{少}{すこ}しだけ\R{見}{み}ました。この\R{章}{しょう}では、その\R{曲線}{きょく|せん}を\R{自分}{じ|ぶん}の\R{手}{て}でかいて、なぜそういう\R{形}{かたち}になるのかを、1つずつ\R{確}{たし}かめます。

%% ---------------- やってみよう・思い出そう ----------------
\begin{yattemiyou}[\R{注射器}{ちゅう|しゃ|き}の\R{空気}{くう|き}をおしちぢめよう]
\noindent\begin{minipage}[c]{0.58\linewidth}
\begin{enumerate}[label={\bkcircled[bkGreen]{\arabic*}}, leftmargin=2.4\zw, itemsep=1pt, topsep=0pt]
\item \R{注射器}{ちゅう|しゃ|き}に\R{空気}{くう|き}を\qty{20}{mL}\R{入}{い}れ、\R{先}{さき}を\R{圧力計}{あつ|りょく|けい}（\R{空気}{くう|き}がおす\R{強}{つよ}さをはかる\R{計器}{けい|き}）につなぐ。
\item ピストンをゆっくりおして、\R{体積}{たい|せき}を16、12、10、8、\qty{5}{mL}にし、そのたびに\R{圧力計}{あつ|りょく|けい}を\R{読}{よ}む。
\item \R{体積}{たい|せき}と\R{圧力}{あつ|りょく}をかけ\R{算}{ざん}してみよう。
\end{enumerate}
\end{minipage}\hfill
\begin{minipage}[c]{0.39\linewidth}\centering
\begin{tikzpicture}[x=0.55mm, y=0.55mm]
  \bkthrhsyringe{(12,0)}{24}
  % 圧力計（丸いメーター）  % noruby
  \draw[bkGray!80, line width=1.2pt] (5,5) -- (-1,5) -- (-1,13);
  \filldraw[fill=white, draw=bkInk, line width=0.9pt] (-1,20) circle (7);
  \foreach \t in {210,180,...,-30} {\draw[bkInk, line width=0.5pt] ($(-1,20)+(\t:5.6)$) -- ($(-1,20)+(\t:6.6)$);}
  \draw[bkRed, line width=1pt] (-1,20) -- ++(60:5);
  \fill[bkInk] (-1,20) circle (0.8);
  \node[font=\sffamily\scriptsize, anchor=south] at (-1,27.5) {kPa};
\end{tikzpicture}
\end{minipage}
\par\vspace{4pt}
\R{学校}{がっ|こう}ではかった\R{結果}{けっ|か}です（\R{圧力}{あつ|りょく}の\R{単位}{たん|い} kPa は\R{次}{つぎ}のページで\R{説明}{せつ|めい}します）。
\begin{bkcenter}
\begin{bktable}{colspec={Q[l,m,wd=33mm] *{6}{Q[c,m,wd=10mm]}}, rowsep=1.5pt}
\R{体積}{たい|せき} $V$（mL） & 20 & 16 & 12 & 10 & 8 & 5 \\
\R{圧力}{あつ|りょく} $p$（kPa） & 101 & 125 & 168 & 201 & 249 & 402 \\
\end{bktable}
\end{bkcenter}
\R{体積}{たい|せき}が\R{半分}{はん|ぶん}の\qty{10}{mL}になると、\R{圧力}{あつ|りょく}はどうなったかな？
\end{yattemiyou}

\noindent\textbf{\R{圧力}{あつ|りょく}とkPa}\par\noindent
\R{空気}{くう|き}がまわりをおす\R{強}{つよ}さを、\qty{1}{m^2}あたりの\R{力}{ちから}で\R{表}{あらわ}したものを\textbf{\R{圧力}{あつ|りょく}}といいます（\R{中}{ちゅう}1の\R{理科}{り|か}）。\bkidx{あつりょく}{\R{圧力}{あつ|りょく}}\R{単位}{たん|い}は Pa（パスカル）で、\qty{1}{m^2}あたり\qty{1}{N}の\R{力}{ちから}が\qty{1}{Pa}です。1000 Pa を1 kPa（キロパスカル）と\R{書}{か}きます（2-10の「k は1000」）。

わたしたちのまわりの\R{空気}{くう|き}の\R{圧力}{あつ|りょく}（\R{大気圧}{たい|き|あつ}）は、\R{約}{やく}\qty{101}{kPa}です。\R{天気予報}{てん|き|よ|ほう}の「1013 hPa（ヘクトパスカル）」と\R{同}{おな}じ\R{大}{おお}きさです（h は100なので、$1013\times100=101300$ Pa）。だから、ピストンをおす\R{前}{まえ}の\R{注射器}{ちゅう|しゃ|き}の\R{中}{なか}も、\R{約}{やく}\qty{101}{kPa}でした。

やってみようの\R{表}{ひょう}で、\R{体積}{たい|せき}が20から10へ\R{半分}{はん|ぶん}になると、\R{圧力}{あつ|りょく}は101から201へ、\R{約}{やく}2\R{倍}{ばい}になりました。20から5へ $\dfrac{1}{4}$ になると、101から402へ\R{約}{やく}4\R{倍}{ばい}です。2-15で\R{学}{まな}んだ\R{反比例}{はん|ぴ|れい}に\R{似}{に}ています。

\begin{omoidasou}
\bkrecall{2-15}{\R{反比例}{はん|ぴ|れい} $y=\dfrac{a}{x}$：$x$ が2\R{倍}{ばい}、3\R{倍}{ばい}になると、$y$ は $\dfrac{1}{2}$ \R{倍}{ばい}、$\dfrac{1}{3}$ \R{倍}{ばい}。$xy$ がいつも\R{同}{おな}じ\R{数}{かず} $a$。}
\bkrecall{1-14}{0でわることはできない。$12\div0$ は\R{計算}{けい|さん}できない。}
\bkrecall{3-2}{\R{点}{てん} $(x,\ y)$ は、\R{原点}{げん|てん}から\R{右}{みぎ}へ $x$、\R{上}{うえ}へ $y$。マイナスなら\R{左}{ひだり}へ・\R{下}{した}へ。\R{座標}{ざ|ひょう}\R{平面}{へい|めん}は4つの\R{区画}{く|かく}に\R{分}{わ}かれる。}
\bkrecall{3-5}{\R{傾}{かたむ}き（\R{変化}{へん|か}の\R{割合}{わり|あい}）＝たての\R{変化量}{へん|か|りょう}÷\R{横}{よこ}の\R{変化量}{へん|か|りょう}。1\R{次関数}{じ|かん|すう} $y=ax+b$ では、いつも $a$。}
\end{omoidasou}

%% ---------------- 数学の道具 ----------------
\bksection{doguu}{\R{数学}{すう|がく}の\R{道具}{どう|ぐ}}

\bksubsection{1}{\R{反比例}{はん|ぴ|れい}の\R{点}{てん}を、\R{自分}{じ|ぶん}でかいてみよう}

\leavevmode\label{bk:3-7:ten}まず、\R{数}{かず}のかんたんな $y=\dfrac{12}{x}$ で\R{考}{かんが}えます。「\R{面積}{めん|せき}が\qty{12}{cm^2}の\R{長方形}{ちょう|ほう|けい}の、\R{横}{よこ} $x$〔cm〕とたて $y$〔cm〕」と\R{同}{おな}じ\R{関係}{かん|けい}です（たて＝\R{面積}{めん|せき}÷\R{横}{よこ}）。12は、1、2、3、4、6、12でわりきれるので、\R{表}{ひょう}がきれいな\R{数}{かず}になります。

\begin{bkcenter}
\begin{bktable}{colspec={Q[l,m,wd=18mm] *{6}{Q[c,m,wd=11mm]}}, rowsep=1.5pt}
$x$ & 1 & 2 & 3 & 4 & 6 & 12 \\
$y=\dfrac{12}{x}$ & 12 & 6 & 4 & 3 & 2 & 1 \\
\SetCell{bg=bkYellow!30} $x\times y$ & 12 & 12 & 12 & 12 & 12 & 12 \\
\end{bktable}
\bktabcaption{$x$ に1つずつ\R{数}{かず}を\R{入}{い}れて、$12\div x$ を\R{計算}{けい|さん}した。}
\end{bkcenter}

\bkthrhstep{1}{\R{表}{ひょう}の\R{組}{くみ}を\R{点}{てん}にする}
$x=2$、$y=6$ の\R{組}{くみ}は、\R{点}{てん} $(2,\ 6)$ です。\R{原点}{げん|てん}から\R{右}{みぎ}へ2、そこから\R{上}{うえ}へ6\R{進}{すす}んだ\R{所}{ところ}にかきます（3-2）。ほかの5つの\R{組}{くみ}も\R{同}{おな}じようにかくと、\R{図}{ず}\ref{bk:3-7:fig-oresen}の6つの\R{黒}{くろ}い\R{点}{てん}になります。

\bkthrhstep{2}{\R{点}{てん}と\R{点}{てん}のあいだは、どうなっている？}
\R{点}{てん}を\R{定規}{じょう|ぎ}でまっすぐつなぐ（おれ\R{線}{せん}にする）と、よいでしょうか。\R{点}{てん}のあいだの $x$ でも、ちゃんと\R{計算}{けい|さん}してみます。$x=1.5$ なら $12\div1.5=8$、$x=5$ なら $12\div5=2.4$ です。

\noindent\begin{minipage}[t]{0.49\linewidth}\centering\vspace{0pt}
\begin{tikzpicture}
\begin{axis}[bkplot, width=60mm, height=56mm, xmin=0, xmax=13.5, ymin=0, ymax=13.5, xtick={1,2,...,12}, ytick={1,2,...,12},
  xticklabels={1,2,3,4,,6,,,,,,12}, yticklabels={1,2,3,4,,6,,8,,,,12}, xlabel={$x$}, ylabel={$y$}]
  \addplot[bkGray, line width=0.9pt, dashed] coordinates {(1,12) (2,6) (3,4) (4,3) (6,2) (12,1)};
  \addplot[only marks, mark=*, mark size=1.9pt, bkInk] coordinates {(1,12) (2,6) (3,4) (4,3) (6,2) (12,1)};
  \addplot[only marks, mark=*, mark size=2.1pt, bkRed] coordinates {(1.5,8) (5,2.4)};
  \node[font=\sffamily\scriptsize, text=bkRed!80!black, anchor=west] at (axis cs:1.75,8.3) {$(1.5,\ 8)$};
  \node[font=\sffamily\scriptsize, text=bkGray, anchor=west] at (axis cs:1.75,10.2) {おれ\R{線}{せん}なら9};
  \node[font=\sffamily\scriptsize, text=bkRed!80!black, anchor=south] at (axis cs:6.4,3.6) {$(5,\ 2.4)$};
\end{axis}
\end{tikzpicture}
\bkfigcaption{おれ\R{線}{せん}（\R{点線}{てん|せん}）と、\R{本当}{ほん|とう}の\R{点}{てん}（\R{赤}{あか}）}\label{bk:3-7:fig-oresen}
\end{minipage}\hfill
\begin{minipage}[t]{0.49\linewidth}\centering\vspace{0pt}
\begin{tikzpicture}
\begin{axis}[bkplot, width=60mm, height=56mm, xmin=0, xmax=13.5, ymin=0, ymax=13.5, xtick={1,2,...,12}, ytick={1,2,...,12},
  xticklabels={1,2,3,4,,6,,,,,,12}, yticklabels={1,2,3,4,,6,,8,,,,12}, xlabel={$x$}, ylabel={$y$}]
  \addplot[bkcurve, domain=0.9:13.3, samples=120] {12/x};
  \addplot[only marks, mark=*, mark size=1.9pt, bkInk] coordinates {(1,12) (2,6) (3,4) (4,3) (6,2) (12,1) (1.5,8) (5,2.4)};
  \node[font=\sffamily\footnotesize, text=bkBlue!80!black, anchor=west] at (axis cs:4,6) {$y=\dfrac{12}{x}$};
\end{axis}
\end{tikzpicture}
\bkfigcaption{\R{点}{てん}を、なめらかな\R{曲線}{きょく|せん}でつなぐ}\label{bk:3-7:fig-kyokusen}
\end{minipage}

\vspace{3pt}
おれ\R{線}{せん}なら、$x=1.5$ では12と6のまん\R{中}{なか}の9になるはずです。でも\R{本当}{ほん|とう}は8で、おれ\R{線}{せん}より\R{下}{した}にあります。$x=5$ でも、おれ\R{線}{せん}の2.5より\R{下}{した}の2.4です。\R{点}{てん}をふやすほど、\R{点}{てん}はなめらかにまがった\R{線}{せん}の\R{上}{うえ}にならびます（\R{図}{ず}\ref{bk:3-7:fig-kyokusen}）。だから\R{反比例}{はん|ぴ|れい}のグラフは、おれ\R{線}{せん}ではなく、\textbf{なめらかな\R{曲線}{きょく|せん}}でつなぎます（3-4の「なめらかな\R{線}{せん}を\R{引}{ひ}く」と\R{同}{おな}じ）。

\bkthrhstep{3}{\R{直線}{ちょく|せん}ではないことを、\R{傾}{かたむ}きで\R{確}{たし}かめる}
もし\R{直線}{ちょく|せん}なら、どの2\R{点}{てん}を\R{選}{えら}んでも\R{傾}{かたむ}きは\R{同}{おな}じです（3-5）。\R{点}{てん} $(1,\ 12)$ から $(2,\ 6)$ までの\R{傾}{かたむ}きと、$(2,\ 6)$ から $(4,\ 3)$ までの\R{傾}{かたむ}きをくらべます。
\[ \frac{6-12}{2-1}=\frac{-6}{1}=-6,\qquad \frac{3-6}{4-2}=\frac{-3}{2}=-1.5 \]
\R{傾}{かたむ}きがちがうので、3つの\R{点}{てん}は1\R{本}{ぽん}の\R{直線}{ちょく|せん}にのっていません。\R{比例}{ひ|れい}のグラフ（2-15）とちがい、\R{反比例}{はん|ぴ|れい}のグラフは\R{直線}{ちょく|せん}ではありません。

\bksubsection{2}{\R{軸}{じく}に\R{近}{ちか}づくけれど、くっつかない}

\leavevmode\label{bk:3-7:jiku}グラフの\R{右}{みぎ}のはしと\R{上}{うえ}のはしは、どうなっているでしょう。$x$ をうんと\R{大}{おお}きくしたり、うんと\R{小}{ちい}さくしたりして\R{計算}{けい|さん}します。

\begin{bkcenter}
\begin{bktable}{colspec={Q[l,m,wd=13mm] *{4}{Q[c,m,wd=10mm]} Q[c,m,wd=5mm] *{3}{Q[c,m,wd=10mm]}}, rowsep=1.5pt, colsep=3pt, cells={font=\small}}
$x$ & 0.01 & 0.1 & 0.5 & 1 & … & 12 & 120 & 1200 \\
$y=\dfrac{12}{x}$ & 1200 & 120 & 24 & 12 & … & 1 & 0.1 & 0.01 \\
\end{bktable}
\bktabcaption{$x$ を\R{小}{ちい}さくすると $y$ はどんどん\R{大}{おお}きく、$x$ を\R{大}{おお}きくすると $y$ はどんどん0に\R{近}{ちか}づく。}
\end{bkcenter}

$x$ を\R{大}{おお}きくするほど、$y$ は0に\R{近}{ちか}づき、グラフは\R{横}{よこ}の\R{軸}{じく}（$x$ \R{軸}{じく}）にどんどん\R{近}{ちか}づきます。$x$ を0に\R{近}{ちか}づけるほど、$y$ はどんどん\R{大}{おお}きくなり、グラフはたての\R{軸}{じく}（$y$ \R{軸}{じく}）に\R{近}{ちか}づきます。では、いつか\R{軸}{じく}にくっつくのでしょうか。2\R{通}{とお}りの\R{方法}{ほう|ほう}で\R{考}{かんが}えます。

\bkthrhhead{\R{考}{かんが}え\R{方}{かた}A：かけ\R{算}{ざん}で\R{考}{かんが}える}
グラフの\R{上}{うえ}の\R{点}{てん}は、どれも $x\times y=12$ をみたします。もしグラフが $x$ \R{軸}{じく}にくっついたら、その\R{点}{てん}では $y=0$ です。すると $x\times0=0$ になり、12になりません。どんな\R{数}{かず}に0をかけても0だからです。だから、$y$ がちょうど0になる\R{点}{てん}は、グラフの\R{上}{うえ}にありません。$y$ \R{軸}{じく}も\R{同}{おな}じで、$x=0$ なら $0\times y=0$ です。そもそも、$x=0$ のときの $y$ は $12\div0$ で、\R{計算}{けい|さん}できません（1-14）。

\bkthrhhead{\R{考}{かんが}え\R{方}{かた}B：\R{長方形}{ちょう|ほう|けい}の\R{面積}{めん|せき}で\R{考}{かんが}える}
\noindent\begin{minipage}[c]{0.5\linewidth}
グラフの\R{上}{うえ}の\R{点}{てん} $(x,\ y)$ と\R{原点}{げん|てん}を、\R{向}{む}かい\R{合}{あ}う\R{角}{かど}にした\R{長方形}{ちょう|ほう|けい}をかきます。\R{横}{よこ}が $x$、たてが $y$ なので、\R{面積}{めん|せき}は $xy=12$ です。

\R{右}{みぎ}の\R{図}{ず}の3つの\R{長方形}{ちょう|ほう|けい}は、どれも\R{面積}{めん|せき}が12です。\R{点}{てん}が\R{右}{みぎ}へ\R{行}{い}くほど、\R{長方形}{ちょう|ほう|けい}は\R{横}{よこ}に\R{長}{なが}く、うすくなります。でも、\R{面積}{めん|せき}を12のままにするには、\R{高}{たか}さを0にすることはできません。\R{高}{たか}さが0の\R{長方形}{ちょう|ほう|けい}は、\R{面積}{めん|せき}が0だからです。
\end{minipage}\hfill
\begin{minipage}[c]{0.47\linewidth}\centering
\begin{tikzpicture}
\begin{axis}[bkplot, width=62mm, height=56mm, xmin=0, xmax=13.5, ymin=0, ymax=13.5, xtick={1,2,3,4,6,12}, ytick={1,2,3,4,6,12}, xlabel={$x$}, ylabel={$y$}, grid=none]
  \fill[bkGreen!22] (axis cs:0,0) rectangle (axis cs:12,1);
  \fill[bkOrange!25] (axis cs:0,0) rectangle (axis cs:4,3);
  \fill[bkBlue!18] (axis cs:0,0) rectangle (axis cs:1,12);
  \draw[bkBlue!80!black, line width=0.7pt] (axis cs:0,0) rectangle (axis cs:1,12);
  \draw[bkOrange!80!black, line width=0.7pt] (axis cs:0,0) rectangle (axis cs:4,3);
  \draw[bkGreen!60!black, line width=0.7pt] (axis cs:0,0) rectangle (axis cs:12,1);
  \addplot[bkcurve, domain=0.9:13.3, samples=120] {12/x};
  \addplot[only marks, mark=*, mark size=1.7pt, bkInk] coordinates {(1,12) (4,3) (12,1)};
  \node[font=\sffamily\scriptsize, text=bkBlue!80!black, anchor=west] at (axis cs:1.3,11) {$1\times12$};
  \node[font=\sffamily\scriptsize, text=bkOrange!80!black, anchor=south west] at (axis cs:4.1,3.1) {$4\times3$};
  \node[font=\sffamily\scriptsize, text=bkGreen!50!black, anchor=south] at (axis cs:9.5,1.3) {$12\times1$};
\end{axis}
\end{tikzpicture}
\bkfigcaption{\R{面積}{めん|せき}はどれも12}
\end{minipage}

\vspace{3pt}
2\R{通}{とお}りのどちらで\R{考}{かんが}えても、グラフは $x$ \R{軸}{じく}にも $y$ \R{軸}{じく}にも、\textbf{\R{近}{ちか}づくけれど、くっつかない}ことがわかりました。\R{原点}{げん|てん} $(0,\ 0)$ も\R{通}{とお}りません。

2-15では、「\R{曲線}{きょく|せん}は、\R{横}{よこ}の\R{軸}{じく}に\R{近}{ちか}づくけれど、くっつきません。\R{比例}{ひ|れい}とちがい、\R{直線}{ちょく|せん}ではなく、\R{原点}{げん|てん}も\R{通}{とお}りません（くわしくは3-7）」と\R{書}{か}きました。そのわけが、これでそろいました。
\begin{itemize}[label={\color{bkOrange}●}, leftmargin=1.6\zw, itemsep=1pt, topsep=2pt]
\item \R{直線}{ちょく|せん}ではない……\R{傾}{かたむ}きが\R{場所}{ば|しょ}によってちがうから（\R{道具}{どう|ぐ}1の③）。
\item \R{軸}{じく}にくっつかない……$xy=12$ なので、$x$ も $y$ も0になれないから（\R{考}{かんが}え\R{方}{かた}A・B）。
\item \R{原点}{げん|てん}を\R{通}{とお}らない……$x=0$ のときの $y$ がないから。
\end{itemize}

\begin{seikaku}
\textbf{「\R{近}{ちか}づくけれど、くっつかない」とは}\par
$y$ を0.001より\R{小}{ちい}さくしたければ $x$ を12000より\R{大}{おお}きく、0.000001より\R{小}{ちい}さくしたければ $x$ を12000000より\R{大}{おお}きくすればよい、というように、どんなに\R{小}{ちい}さい\R{数}{かず}を\R{決}{き}めても、$y$ をそれより\R{小}{ちい}さくできます。それでも $y$ は0にはなりません。グラフがこのように\R{近}{ちか}づいていく\R{直線}{ちょく|せん}を、\R{高校}{こう|こう}の\R{数学}{すう|がく}では\textbf{\R{漸近線}{ぜん|きん|せん}}とよびます。$y=\dfrac{12}{x}$ では、$x$ \R{軸}{じく}と $y$ \R{軸}{じく}が\R{漸近線}{ぜん|きん|せん}です。
\end{seikaku}

\bksubsection{3}{$x$ が\R{負}{ふ}のがわ：2\R{本}{ほん}の\R{枝}{えだ}}

\leavevmode\label{bk:3-7:eda}\R{面積}{めん|せき}や\R{体積}{たい|せき}はマイナスになりませんが、\R{式}{しき} $y=\dfrac{12}{x}$ には、マイナスの\R{数}{かず}も\R{入}{い}れられます。$x=-2$ なら、$y=12\div(-2)=-6$ です（プラス÷マイナスはマイナス、2-2）。

\begin{bkcenter}
\begin{bktable}{colspec={Q[l,m,wd=16mm] *{6}{Q[c,m,wd=11mm]}}, rowsep=1.5pt}
$x$ & $-12$ & $-6$ & $-4$ & $-3$ & $-2$ & $-1$ \\
$y=\dfrac{12}{x}$ & $-1$ & $-2$ & $-3$ & $-4$ & $-6$ & $-12$ \\
\SetCell{bg=bkYellow!30} $x\times y$ & 12 & 12 & 12 & 12 & 12 & 12 \\
\end{bktable}
\bktabcaption{マイナス×マイナスはプラスなので、$x\times y$ はやはり12。}
\end{bkcenter}

\noindent\begin{minipage}[c]{0.5\linewidth}
これらの\R{点}{てん}は、$x$ も $y$ もマイナスなので、\R{左下}{ひだり|した}の\R{区画}{く|かく}（3-2）に\R{入}{はい}ります。なめらかにつなぐと、\R{右上}{みぎ|うえ}の\R{曲線}{きょく|せん}とそっくりな\R{形}{かたち}の\R{曲線}{きょく|せん}が、もう1\R{本}{ぽん}できます。

\R{点}{てん} $(2,\ 6)$ と $(-2,\ -6)$ をくらべてみましょう。どちらも\R{原点}{げん|てん}から\R{同}{おな}じ\R{距離}{きょ|り}にあり、\R{原点}{げん|てん}をはさんでちょうど\R{反対}{はん|たい}がわです（\R{図}{ず}の\R{点線}{てん|せん}）。どの\R{点}{てん}でも\R{同}{おな}じなので、2\R{本}{ほん}の\R{曲線}{きょく|せん}は、\R{原点}{げん|てん}を\R{中心}{ちゅう|しん}に\R{半回転}{はん|かい|てん}させると\R{重}{かさ}なります。
\end{minipage}\hfill
\begin{minipage}[c]{0.47\linewidth}\centering
\begin{tikzpicture}
\begin{axis}[bkplot, width=62mm, height=62mm, xmin=-13.5, xmax=13.5, ymin=-13.5, ymax=13.5, xtick={-12,-6,6,12}, ytick={-12,-6,6,12}, xlabel={$x$}, ylabel={$y$}]
  \addplot[bkcurve, domain=0.9:13.3, samples=120] {12/x};
  \addplot[bkcurve, domain=-13.3:-0.9, samples=120] {12/x};
  \draw[bkGray, line width=0.7pt, dash pattern=on 2pt off 1.5pt] (axis cs:-2,-6) -- (axis cs:2,6);
  \addplot[only marks, mark=*, mark size=1.8pt, bkRed] coordinates {(2,6) (-2,-6)};
  \node[font=\sffamily\scriptsize, anchor=west, text=bkRed!80!black] at (axis cs:2.4,6.6) {$(2,\ 6)$};
  \node[font=\sffamily\scriptsize, anchor=east, text=bkRed!80!black] at (axis cs:-2.4,-6.6) {$(-2,\ -6)$};
\end{axis}
\end{tikzpicture}
\bkfigcaption{$y=\dfrac{12}{x}$ の2\R{本}{ほん}の\R{枝}{えだ}}
\end{minipage}

\vspace{3pt}
\R{反比例}{はん|ぴ|れい}のグラフは、このように\R{離}{はな}れた2\R{本}{ほん}の\R{曲線}{きょく|せん}でできています。2\R{本}{ほん}は $x=0$ のところでつながりません。$x=0$ のときの $y$ がないからです。この\R{曲線}{きょく|せん}を\textbf{\R{双曲線}{そう|きょく|せん}}といい、1\R{本}{ぽん}1\R{本}{ぽん}を\R{双曲線}{そう|きょく|せん}の\textbf{\R{枝}{えだ}}といいます。\bkidx{そうきょくせん}{\R{双曲線}{そう|きょく|せん}}「\R{双}{そう}」は「2つで1\R{組}{くみ}」という\R{意味}{い|み}です。

\R{体積}{たい|せき}や\R{速}{はや}さのように、マイナスにならない\R{量}{りょう}の\R{反比例}{はん|ぴ|れい}では、$x>0$ の\R{範囲}{はん|い}（\R{変域}{へん|いき}、3-1）だけを\R{考}{かんが}えるので、\R{右上}{みぎ|うえ}の\R{枝}{えだ}だけを\R{使}{つか}います。

\bksubsection{4}{$a$ を\R{変}{か}えると、グラフはどう\R{変}{か}わる？}

\leavevmode\label{bk:3-7:fu}\bkthrhhead{$a$ が\R{大}{おお}きいと、\R{原点}{げん|てん}から\R{遠}{とお}くなる}
$y=\dfrac{4}{x}$、$y=\dfrac{12}{x}$、$y=\dfrac{24}{x}$ をくらべます。$x=2$ のときの $y$ は、それぞれ2、6、12です。$a$ が\R{大}{おお}きいほど、\R{同}{おな}じ $x$ での $y$ が\R{大}{おお}きく、\R{曲線}{きょく|せん}は\R{原点}{げん|てん}から\R{遠}{とお}い\R{所}{ところ}を\R{通}{とお}ります（\R{図}{ず}\ref{bk:3-7:fig-a}）。\R{面積}{めん|せき}で\R{考}{かんが}えると、\R{面積}{めん|せき} $a$ の\R{長方形}{ちょう|ほう|けい}の\R{角}{かど}をつないだ\R{線}{せん}なので、\R{面積}{めん|せき}が\R{大}{おお}きいほど\R{外}{そと}がわになるのです。

\bkthrhhead{$a$ が\R{負}{ふ}のとき}
$y=-\dfrac{12}{x}$ の\R{表}{ひょう}を\R{作}{つく}ります。$a=-12$ です。

\begin{bkcenter}
\begin{bktable}{colspec={Q[l,m,wd=18mm] *{8}{Q[c,m,wd=9.5mm]}}, rowsep=1.5pt, cells={font=\small}}
$x$ & $-12$ & $-4$ & $-2$ & $-1$ & 1 & 2 & 4 & 12 \\
$y=-\dfrac{12}{x}$ & 1 & 3 & 6 & 12 & $-12$ & $-6$ & $-3$ & $-1$ \\
\end{bktable}
\bktabcaption{$y=\dfrac{12}{x}$ の $y$ の\R{符号}{ふ|ごう}が、すべて\R{反対}{はん|たい}になっている。}
\end{bkcenter}

\noindent\begin{minipage}[t]{0.49\linewidth}\centering\vspace{0pt}
\begin{tikzpicture}
\begin{axis}[bkplot, width=60mm, height=56mm, xmin=0, xmax=13.5, ymin=0, ymax=13.5, xtick={2,4,6,8,10,12}, ytick={2,4,...,12}, xlabel={$x$}, ylabel={$y$}]
  \addplot[bkcurve3, domain=0.31:13.3, samples=120] {4/x};
  \addplot[bkcurve, domain=0.9:13.3, samples=120] {12/x};
  \addplot[bkRed, line width=1.7pt, domain=1.8:13.3, samples=120] {24/x};
  \node[font=\sffamily\scriptsize, text=bkGreen!50!black, anchor=south west, inner sep=1pt] at (axis cs:1.7,2.4) {$a=4$};
  \node[font=\sffamily\scriptsize, text=bkBlue!80!black, anchor=south west, inner sep=1pt] at (axis cs:3.6,3.6) {$a=12$};
  \node[font=\sffamily\scriptsize, text=bkRed!80!black, anchor=south west, inner sep=1pt] at (axis cs:5,5.2) {$a=24$};
\end{axis}
\end{tikzpicture}
\bkfigcaption{$a$ が\R{大}{おお}きいほど\R{外}{そと}がわ}\label{bk:3-7:fig-a}
\end{minipage}\hfill
\begin{minipage}[t]{0.49\linewidth}\centering\vspace{0pt}
\begin{tikzpicture}
\begin{axis}[bkplot, width=60mm, height=56mm, xmin=-13.5, xmax=13.5, ymin=-13.5, ymax=13.5, xtick={-12,-6,6,12}, ytick={-12,-6,6,12}, xlabel={$x$}, ylabel={$y$}]
  \addplot[bkBlue!35, line width=1.1pt, dashed, domain=0.9:13.3, samples=100] {12/x};
  \addplot[bkBlue!35, line width=1.1pt, dashed, domain=-13.3:-0.9, samples=100] {12/x};
  \addplot[bkcurve2, domain=0.9:13.3, samples=120] {-12/x};
  \addplot[bkcurve2, domain=-13.3:-0.9, samples=120] {-12/x};
  \node[font=\sffamily\scriptsize, text=bkOrange!80!black, anchor=north west, inner sep=1pt] at (axis cs:2.6,-5) {$y=-\dfrac{12}{x}$};
  \node[font=\sffamily\scriptsize, text=bkBlue!70!black, anchor=south west, inner sep=1pt] at (axis cs:3,5) {$y=\dfrac{12}{x}$};
\end{axis}
\end{tikzpicture}
\bkfigcaption{$a$ が\R{負}{ふ}なら\R{右下}{みぎ|した}と\R{左上}{ひだり|うえ}}\label{bk:3-7:fig-fu}
\end{minipage}

\vspace{3pt}
$y=-\dfrac{12}{x}$ のグラフ（\R{図}{ず}\ref{bk:3-7:fig-fu}のオレンジ）は、$y=\dfrac{12}{x}$ のグラフ（\R{点線}{てん|せん}）を、$x$ \R{軸}{じく}で\R{折}{お}り\R{返}{かえ}した\R{形}{かたち}です。\R{同}{おな}じ $x$ で、$y$ の\R{符号}{ふ|ごう}だけが\R{反対}{はん|たい}になるからです。\R{枝}{えだ}は\R{右下}{みぎ|した}と\R{左上}{ひだり|うえ}の\R{区画}{く|かく}に\R{入}{はい}ります。$xy=-12$ で、$x$ と $y$ の\R{符号}{ふ|ごう}がいつもちがうからです。

ここで\R{気}{き}をつけたいことがあります。$x>0$ のがわでは、$x$ が1、2、4とふえると、$y$ は $-12$、$-6$、$-3$ と\textbf{ふえて}います。グラフは\R{右上}{みぎ|うえ}がりです。「\R{反比例}{はん|ぴ|れい}は、ふえるとへる」とはかぎりません。2-15の「$x$ が2\R{倍}{ばい}になると $y$ は $\dfrac{1}{2}$ \R{倍}{ばい}」は、$a<0$ でもそのまま\R{成}{な}り\R{立}{た}ちます（$-12$ の $\dfrac{1}{2}$ \R{倍}{ばい}は $-6$）。$\dfrac{1}{2}$ \R{倍}{ばい}になると0に\R{近}{ちか}づくので、マイナスの\R{数}{かず}は\R{大}{おお}きくなるのです。

\bksubsection{5}{\R{変化}{へん|か}の\R{割合}{わり|あい}は\R{一定}{いっ|てい}ではない}

\leavevmode\label{bk:3-7:wariai}3-5では、1\R{次関数}{じ|かん|すう} $y=ax+b$ の\R{変化}{へん|か}の\R{割合}{わり|あい}は、どこでもいつも $a$ でした。\R{反比例}{はん|ぴ|れい}ではどうでしょう。$y=\dfrac{12}{x}$ で、$x$ の\R{区間}{く|かん}ごとに、\R{変化}{へん|か}の\R{割合}{わり|あい}＝（$y$ の\R{変化量}{へん|か|りょう}）÷（$x$ の\R{変化量}{へん|か|りょう}）を\R{計算}{けい|さん}します。

\begin{bkcenter}
\begin{bktable}{colspec={Q[c,m,wd=24mm] Q[c,m,wd=26mm] Q[c,m,wd=26mm] Q[c,m,wd=28mm]}, rowsep=1.5pt, cells={font=\small}}
$x$ の\R{区間}{く|かん} & $x$ の\R{変化量}{へん|か|りょう} & $y$ の\R{変化量}{へん|か|りょう} & \R{変化}{へん|か}の\R{割合}{わり|あい} \\
1 → 2 & 1 & $6-12=-6$ & $-6\div1=-6$ \\
2 → 3 & 1 & $4-6=-2$ & $-2\div1=-2$ \\
3 → 4 & 1 & $3-4=-1$ & $-1\div1=-1$ \\
4 → 6 & 2 & $2-3=-1$ & $-1\div2=-0.5$ \\
6 → 12 & 6 & $1-2=-1$ & $-1\div6=-\dfrac{1}{6}$ \\
\end{bktable}
\bktabcaption{$y=\dfrac{12}{x}$ の\R{変化}{へん|か}の\R{割合}{わり|あい}。\R{右}{みぎ}へ\R{行}{い}くほど0に\R{近}{ちか}づく。}
\end{bkcenter}

\noindent\begin{minipage}[c]{0.5\linewidth}
\R{変化}{へん|か}の\R{割合}{わり|あい}は\R{区間}{く|かん}ごとにちがい、\R{一定}{いっ|てい}ではありません。$x$ を\R{同}{おな}じ1ずつふやしても、はじめは $y$ が6もへり、\R{次}{つぎ}は2、その\R{次}{つぎ}は1しかへりません。\R{右}{みぎ}の\R{図}{ず}の\R{階段}{かい|だん}（\R{緑}{みどり}）の\R{段}{だん}の\R{高}{たか}さが、だんだん\R{低}{ひく}くなっています。
\end{minipage}\hfill
\begin{minipage}[c]{0.47\linewidth}\centering
\begin{tikzpicture}
\begin{axis}[bkplot, width=62mm, height=58mm, xmin=0, xmax=13.5, ymin=0, ymax=13.5, xtick={1,2,3,4,6,12}, ytick={1,2,3,4,6,12}, xlabel={$x$}, ylabel={$y$}]
  \addplot[bkGray, line width=0.9pt, dashed, domain=0:13] {-x+13};
  \addplot[bkcurve, domain=0.9:13.3, samples=120] {12/x};
  \addplot[bkGreen!70!black, line width=1.1pt] coordinates {(1,12) (2,12) (2,6) (3,6) (3,4) (4,4) (4,3)};
  \node[font=\sffamily\scriptsize, text=bkGreen!50!black, anchor=west] at (axis cs:2.15,9) {$-6$};
  \node[font=\sffamily\scriptsize, text=bkGreen!50!black, anchor=west] at (axis cs:3.15,5) {$-2$};
  \node[font=\sffamily\scriptsize, text=bkGreen!50!black, anchor=north] at (axis cs:4.2,2.8) {$-1$};
  \node[font=\sffamily\scriptsize, text=bkGray, anchor=south west, inner sep=1pt] at (axis cs:8,5.2) {$y=-x+13$};
\end{axis}
\end{tikzpicture}
\end{minipage}

\vspace{3pt}
\R{同}{おな}じ2\R{点}{てん} $(1,\ 12)$ と $(12,\ 1)$ を\R{通}{とお}る\R{直線}{ちょく|せん} $y=-x+13$（\R{点線}{てん|せん}）は、どこでも\R{変化}{へん|か}の\R{割合}{わり|あい}が $-1$ です。\R{反比例}{はん|ぴ|れい}のグラフは、はじめは\R{急}{きゅう}な\R{下}{くだ}り\R{坂}{ざか}で、\R{右}{みぎ}へ\R{行}{い}くほどゆるやかになる\R{坂}{さか}なのです。

2-15の「$x$ が2\R{倍}{ばい}になると、$y$ は $\dfrac{1}{2}$ \R{倍}{ばい}」は、\R{何倍}{なん|ばい}になるか（わり\R{算}{ざん}の\R{見方}{み|かた}）のきまりでした。3-5の\R{変化}{へん|か}の\R{割合}{わり|あい}は、\R{何}{なん}だけふえるか（ひき\R{算}{ざん}の\R{見方}{み|かた}）です。\R{反比例}{はん|ぴ|れい}は、わり\R{算}{ざん}の\R{見方}{み|かた}ではきれいなきまりがあるのに、ひき\R{算}{ざん}の\R{見方}{み|かた}では\R{一定}{いっ|てい}になりません。

\begin{motto}[3]
\textbf{\R{変化}{へん|か}の\R{割合}{わり|あい}を、1つの\R{式}{しき}で}\par
\R{表}{ひょう}の\R{変化}{へん|か}の\R{割合}{わり|あい}を、\R{区間}{く|かん}の\R{両}{りょう}はしの $x$ で\R{見}{み}てみます。1→2では $-6=-\dfrac{12}{1\times2}$、2→3では $-2=-\dfrac{12}{2\times3}$、4→6では $-0.5=-\dfrac{12}{4\times6}$ です。どれも「$-12\div$（\R{両}{りょう}はしの $x$ の\R{積}{せき}）」になっています。$x$ が $p$ から $q$ まで\R{変}{か}わるとき、
\[ \frac{\dfrac{12}{q}-\dfrac{12}{p}}{q-p}=-\frac{12}{pq} \]
となるのです（\R{分子}{ぶん|し}を\R{通分}{つう|ぶん}すると $\dfrac{12p-12q}{pq}=-\dfrac{12(q-p)}{pq}$ で、これを $q-p$ でわる）。$p$ と $q$ がとても\R{近}{ちか}いと、$pq$ は $x^2$ にほぼ\R{等}{ひと}しくなり、\R{変化}{へん|か}の\R{割合}{わり|あい}は $-\dfrac{12}{x^2}$ に\R{近}{ちか}づきます。これは\R{第}{だい}5\R{巻}{かん}の5-4で、\R{微分}{び|ぶん}を\R{使}{つか}って\R{確}{たし}かめる\R{形}{かたち}です。
\end{motto}

\bksubsection{6}{$xy$ が\R{一定}{いっ|てい}を\R{見}{み}つける：$pV$ が\R{一定}{いっ|てい}}

\leavevmode\label{bk:3-7:pv}やってみようの\R{表}{ひょう}にもどりましょう。2-15のように、\R{体積}{たい|せき} $V$ と\R{圧力}{あつ|りょく} $p$ をかけ\R{算}{ざん}します。

\begin{bkcenter}
\begin{bktable}{colspec={Q[l,m,wd=30mm] *{6}{Q[c,m,wd=11mm]}}, rowsep=1.5pt}
\R{体積}{たい|せき} $V$（mL） & 20 & 16 & 12 & 10 & 8 & 5 \\
\R{圧力}{あつ|りょく} $p$（kPa） & 101 & 125 & 168 & 201 & 249 & 402 \\
\SetCell{bg=bkYellow!30} $p\times V$ & 2020 & 2000 & 2016 & 2010 & 1992 & 2010 \\
\end{bktable}
\bktabcaption{$p\times V$ は、どれも2000に\R{近}{ちか}い。}
\end{bkcenter}

$p\times V$ は、どれも2000の\R{近}{ちか}くです。\R{少}{すこ}しずつちがうのは、はかるときのずれです（2-6）。ならすと $pV=2000$ と\R{考}{かんが}えられます。\R{両辺}{りょう|へん}を $V$ でわると（2-14）、$p=\dfrac{2000}{V}$ です。

\noindent\begin{minipage}[c]{0.47\linewidth}
$p$ は $V$ に\R{反比例}{はん|ぴ|れい}します。$a=2000$ の\R{反比例}{はん|ぴ|れい}です。\R{体積}{たい|せき}はマイナスにならないので、$V>0$ の\R{枝}{えだ}だけをかきます。はかった\R{点}{てん}（\R{黒}{くろ}）は、\R{曲線}{きょく|せん} $p=\dfrac{2000}{V}$ のすぐ\R{近}{ちか}くにならびました。
\end{minipage}\hfill
\begin{minipage}[c]{0.5\linewidth}\centering
\begin{tikzpicture}
\begin{axis}[bkplot, width=56mm, height=56mm, xmin=0, xmax=23, ymin=0, ymax=650, xtick={5,10,15,20}, ytick={100,200,...,600}, xlabel={$V$（mL）}, ylabel={$p$（kPa）}]
  \addplot[bkcurve, domain=3.2:22.5, samples=120] {2000/x};
  \addplot[only marks, mark=*, mark size=1.7pt, bkInk] coordinates {(20,101) (16,125) (12,168) (10,201) (8,249) (5,402)};
  \draw[bkRed, line width=0.7pt, dash pattern=on 2pt off 1.5pt] (axis cs:20,0) -- (axis cs:20,100) (axis cs:10,0) -- (axis cs:10,200) (axis cs:5,0) -- (axis cs:5,400);
  \node[font=\sffamily\scriptsize, text=bkBlue!80!black, anchor=west] at (axis cs:7,470) {$p=\dfrac{2000}{V}$};
\end{axis}
\end{tikzpicture}
\end{minipage}

\vspace{3pt}
\bkhatenaref の\R{答}{こた}え：\R{体積}{たい|せき}を\R{半分}{はん|ぶん}の\qty{10}{mL}にすると、$p=2000\div10=200$ で、おし\R{返}{かえ}す\R{強}{つよ}さ（\R{圧力}{あつ|りょく}）は\R{約}{やく}2\R{倍}{ばい}です。$\dfrac{1}{4}$ の\qty{5}{mL}なら\R{約}{やく}4\R{倍}{ばい}、$\dfrac{1}{10}$ の\qty{2}{mL}なら $2000\div2=1000$ で\R{約}{やく}10\R{倍}{ばい}になります（グラフの\R{赤}{あか}い\R{点線}{てん|せん}）。\R{空気}{くう|き}を0にすることはできません。\R{体積}{たい|せき}を0に\R{近}{ちか}づけるほど\R{圧力}{あつ|りょく}はどこまでも\R{大}{おお}きくなり、グラフは $p$ \R{軸}{じく}にくっつかないからです。グラフの\R{形}{かたち}は、まっすぐな\R{線}{せん}ではなく\R{双曲線}{そう|きょく|せん}です。

\R{温度}{おん|ど}が\R{変}{か}わらなければ、\R{気体}{き|たい}の\R{圧力}{あつ|りょく}と\R{体積}{たい|せき}の\R{積}{せき} $pV$ は\R{一定}{いっ|てい}です。このきまりを\textbf{ボイルの\R{法則}{ほう|そく}}といいます。\bkidx{ぼいるのほうそく}{ボイルの\R{法則}{ほう|そく}}やってみようで「ゆっくりおす」としたのは、\R{速}{はや}くおすと\R{空気}{くう|き}の\R{温度}{おん|ど}が\R{上}{あ}がってしまうからです。

\begin{column}{なぜ、おしちぢめると\R{強}{つよ}くおし\R{返}{かえ}すの？}
\noindent\begin{minipage}[c]{0.66\linewidth}
\R{空気}{くう|き}は、とても\R{小}{ちい}さなつぶ（\R{分子}{ぶん|し}）がたくさん\R{飛}{と}び\R{回}{まわ}ってできています。\R{分子}{ぶん|し}がピストンにぶつかるときの\R{力}{ちから}を\R{全部}{ぜん|ぶ}あわせたものが、\R{空気}{くう|き}がおす\R{力}{ちから}です。\R{体積}{たい|せき}を\R{半分}{はん|ぶん}にすると、\R{同}{おな}じ\R{数}{かず}の\R{分子}{ぶん|し}が\R{半分}{はん|ぶん}の\R{部屋}{へ|や}につめこまれ、1\R{秒}{びょう}あたりにピストンにぶつかる\R{回数}{かい|すう}が2\R{倍}{ばい}になります。だから\R{圧力}{あつ|りょく}も2\R{倍}{ばい}になるのです。
\end{minipage}\hfill
\begin{minipage}[c]{0.3\linewidth}\centering
\begin{tikzpicture}[x=1mm, y=1mm]
  \foreach \yy/\w in {14/28, 0/14} {
    \fill[bkBlue!10] (0,\yy) rectangle (\w,\yy+10);
    \draw[bkGray!90!black, line width=0.8pt] (0,\yy) -- (0,\yy+10) (0,\yy) -- (30,\yy) (0,\yy+10) -- (30,\yy+10);
    \fill[bkInk!70] (\w,\yy+0.3) rectangle (\w+2,\yy+9.7);}
  \foreach \px/\py in {3/17,8/21.5,13/16.5,18/22,23/18,6/20,16/19.5,25/21.5} {\fill[bkBlue!60] (\px,\py) circle (0.7);}
  \foreach \px/\py in {2/3,4.5/7.5,7/2.5,9.5/6.5,12/3.5,3.5/5,8/4.5,11/8} {\fill[bkBlue!60] (\px,\py) circle (0.7);}
\end{tikzpicture}
\end{minipage}
\end{column}

\bkthrhhead{\R{反比例}{はん|ぴ|れい}かどうかを、データから\R{見分}{み|わ}ける\R{手順}{て|じゅん}}
\begin{enumerate}[label={\bkcircled[bkOrange]{\arabic*}}, leftmargin=2.4\zw, itemsep=1pt, topsep=2pt]
\item \R{表}{ひょう}のたての\R{組}{くみ}ごとに、$x\times y$ を\R{計算}{けい|さん}する。
\item どれも\R{同}{おな}じ（はかった\R{値}{あたい}なら、ほぼ\R{同}{おな}じ）なら\R{反比例}{はん|ぴ|れい}。その\R{値}{あたい}が $a$。
\item $y=\dfrac{a}{x}$ と\R{書}{か}き、ほしい $x$ や $y$ を\R{計算}{けい|さん}する。
\end{enumerate}

\begin{bkcenter}
\begin{bktable}{colspec={Q[l,m,wd=26mm] Q[l,m,wd=44mm] Q[l,m,wd=48mm]}, rowsep=1.5pt, cells={font=\small}, row{1}={font=\sffamily\bfseries\small}}
 & \R{比例}{ひ|れい} $y=ax$ & \R{反比例}{はん|ぴ|れい} $y=\dfrac{a}{x}$ \\
\R{一定}{いっ|てい}になるもの & $y\div x=a$ & $x\times y=a$ \\
$x$ を2\R{倍}{ばい}にすると & $y$ も2\R{倍}{ばい} & $y$ は $\dfrac{1}{2}$ \R{倍}{ばい} \\
グラフ & \R{原点}{げん|てん}を\R{通}{とお}る\R{直線}{ちょく|せん} & \R{原点}{げん|てん}を\R{通}{とお}らない\R{双曲線}{そう|きょく|せん} \\
\R{軸}{じく}との\R{関係}{かん|けい} & \R{原点}{げん|てん}で\R{交}{まじ}わる & \R{近}{ちか}づくが、くっつかない \\
\R{変化}{へん|か}の\R{割合}{わり|あい} & いつも $a$ & \R{区間}{く|かん}ごとにちがう \\
\end{bktable}
\bktabcaption{\R{比例}{ひ|れい}と\R{反比例}{はん|ぴ|れい}のくらべ}
\end{bkcenter}

\begin{yattemiyou}[\R{表計算}{ひょう|けい|さん}で $x=0$ を\R{入}{い}れてみよう]
\noindent\begin{minipage}[c]{0.5\linewidth}
\begin{enumerate}[label={\bkcircled[bkGreen]{\arabic*}}, leftmargin=2.4\zw, itemsep=1pt, topsep=0pt]
\item A\R{列}{れつ}に $x$ の\R{値}{あたい}を\R{入}{い}れ、B2に \texttt{=12/A2} と\R{入}{い}れて、\R{下}{した}へコピーする（2-18）。
\item $x=0$ の\R{行}{ぎょう}は、どうなるかな？
\item A\R{列}{れつ}とB\R{列}{れつ}で\R{散布図}{さん|ぷ|ず}をかこう。
\end{enumerate}
\end{minipage}\hfill
\begin{minipage}[c]{0.47\linewidth}\centering
\bksheet{2}{20}{%
1 & $x$ & $y$ \\
2 & $-4$ & $-3$ \\
3 & $-1$ & $-12$ \\
4 & 0 & {\color{bkRed}\bfseries \#DIV/0!} \\
5 & 1 & 12 \\
6 & 4 & 3 \\}
\end{minipage}
\par\vspace{4pt}
\texttt{\#DIV/0!} は「0でわった」（division by 0）というエラーの\R{印}{しるし}です。\R{表計算}{ひょう|けい|さん}も、$x=0$ のときの $y$ は\R{計算}{けい|さん}できません。だから\R{散布図}{さん|ぷ|ず}には、$x=0$ の\R{点}{てん}がかかれず、2\R{本}{ほん}の\R{枝}{えだ}は\R{離}{はな}れたままです。
\end{yattemiyou}

\begin{taisetsu}
\R{反比例}{はん|ぴ|れい} $y=\dfrac{a}{x}$ のグラフは\R{双曲線}{そう|きょく|せん}。\R{直線}{ちょく|せん}ではなく、\R{原点}{げん|てん}も\R{通}{とお}らない。\par\vspace{2pt}
$x$ \R{軸}{じく}・$y$ \R{軸}{じく}に\R{近}{ちか}づくが、くっつかない（$xy=a$ なので、$x$ も $y$ も0になれない）。\par\vspace{2pt}
$a>0$ なら\R{右上}{みぎ|うえ}と\R{左下}{ひだり|した}、$a<0$ なら\R{右下}{みぎ|した}と\R{左上}{ひだり|うえ}に2\R{本}{ほん}の\R{枝}{えだ}。\par\vspace{2pt}
\R{変化}{へん|か}の\R{割合}{わり|あい}は\R{一定}{いっ|てい}ではない。
\end{taisetsu}

\begin{kotoba}
\small
\bkword{\R{双曲線}{そう|きょく|せん}・\R{枝}{えだ}}{\R{反比例}{はん|ぴ|れい}のグラフの\R{形}{かたち}。\R{離}{はな}れた2\R{本}{ほん}の\R{曲線}{きょく|せん}（\R{枝}{えだ}）でできている。}
\bkword{\R{圧力}{あつ|りょく}（Pa）}{\qty{1}{m^2}あたりにおす\R{力}{ちから}。1 kPa＝1000 Pa。まわりの\R{空気}{くう|き}は\R{約}{やく}\qty{101}{kPa}。}
\bkword{ボイルの\R{法則}{ほう|そく}}{\R{温度}{おん|ど}が\R{変}{か}わらないとき、\R{気体}{き|たい}の\R{圧力}{あつ|りょく}と\R{体積}{たい|せき}の\R{積}{せき} $pV$ は\R{一定}{いっ|てい}。}
\bkword{\R{漸近線}{ぜん|きん|せん}}{グラフが\R{近}{ちか}づくけれどくっつかない\R{直線}{ちょく|せん}（\R{高校}{こう|こう}のことば）。}
\end{kotoba}

%% ---------------- 例題 ----------------
\bksection{reidai}{\R{例題}{れい|だい}でたしかめよう}

\begin{reidai}[\R{反比例}{はん|ぴ|れい}の\R{式}{しき}を\R{作}{つく}って\R{使}{つか}う]
$y$ は $x$ に\R{反比例}{はん|ぴ|れい}し、$x=3$ のとき $y=8$ です。（1）$y$ を $x$ の\R{式}{しき}で\R{表}{あらわ}しましょう。（2）$x=6$、$x=-4$ のときの $y$ を\R{求}{もと}めましょう。（3）$y=2$ になる $x$ を\R{求}{もと}めましょう。\par
\kangaekata \R{反比例}{はん|ぴ|れい}なので $y=\dfrac{a}{x}$ の\R{形}{かたち}で、$xy=a$ です。わかっている1\R{組}{くみ}の $x$ と $y$ をかけると、$a$ が\R{求}{もと}まります。\par
\shiki（1）\nobreak $a=xy=3\times8=24$\quad だから $y=\dfrac{24}{x}$\par
\hspace*{2.6\zw}（2）\nobreak $x=6$ を\R{代入}{だい|にゅう}して $y=\dfrac{24}{6}=4$\par
\hspace*{2.6\zw}\phantom{（2）}$x=-4$ を\R{代入}{だい|にゅう}して $y=\dfrac{24}{-4}=-6$\quad（＋÷－は－）\par
\hspace*{2.6\zw}（3）\nobreak $xy=24$ に $y=2$ を\R{入}{い}れて $2x=24$。\R{両辺}{りょう|へん}を2でわって $x=12$\par
\kotae（1）\nobreak $y=\dfrac{24}{x}$\quad（2）\nobreak $x=6$ のとき4、$x=-4$ のとき $-6$\quad（3）\nobreak $x=12$\par
\vspace{2pt}
\noindent\begin{minipage}[c]{0.56\linewidth}
\textbf{グラフで\R{確}{たし}かめる}\par
\R{求}{もと}めた\R{点}{てん} $(3,\ 8)$、$(6,\ 4)$、$(12,\ 2)$、$(-4,\ -6)$ をかきこむと、どれも $y=\dfrac{24}{x}$ の\R{双曲線}{そう|きょく|せん}の\R{上}{うえ}にのっています。どの\R{点}{てん}も、$x\times y$ が24だからです。
\end{minipage}\hfill
\begin{minipage}[c]{0.4\linewidth}\centering
\begin{tikzpicture}
\begin{axis}[bkplot, width=46mm, height=44mm, xmin=-14, xmax=14, ymin=-14, ymax=14, xtick={-12,-6,6,12}, ytick={-12,-6,6,12}, xlabel={$x$}, ylabel={$y$}]
  \addplot[bkcurve, domain=1.75:13.5, samples=100] {24/x};
  \addplot[bkcurve, domain=-13.5:-1.75, samples=100] {24/x};
  \addplot[only marks, mark=*, mark size=1.6pt, bkRed] coordinates {(3,8) (6,4) (12,2) (-4,-6)};
\end{axis}
\end{tikzpicture}
\end{minipage}
\end{reidai}

\begin{jibun}
$y$ は $x$ に\R{反比例}{はん|ぴ|れい}し、$x=4$ のとき $y=-3$ です。（1）$y$ を $x$ の\R{式}{しき}で\R{表}{あらわ}しましょう。（2）$x=2$、$x=-6$ のときの $y$ を\R{求}{もと}めましょう。（3）$y=6$ になる $x$ を\R{求}{もと}めましょう。
\bkjans{（1）\nobreak $a=4\times(-3)=-12$ で $y=-\dfrac{12}{x}$。（2）\nobreak $x=2$ のとき $-6$、$x=-6$ のとき2。（3）\nobreak $6x=-12$ で $x=-2$}
\end{jibun}

\begin{reidai}[\R{表}{ひょう}を\R{作}{つく}って、グラフをかく]
$y=-\dfrac{8}{x}$ について、$x=-8$、$-4$、$-2$、$-1$、1、2、4、8 のときの $y$ の\R{表}{ひょう}を\R{作}{つく}り、グラフをかきましょう。グラフはどの\R{区画}{く|かく}を\R{通}{とお}りますか。\par
\kangaekata $x$ を1つずつ\R{代入}{だい|にゅう}します。マイナスの\R{数}{かず}はかっこをつけて\R{代入}{だい|にゅう}します（2-13）。\R{点}{てん}は\R{直線}{ちょく|せん}でつながず、なめらかな\R{曲線}{きょく|せん}でつなぎます。$x=0$ の\R{所}{ところ}ではつなぎません。\par
\shiki $x=-8$ なら $y=-\dfrac{8}{(-8)}=1$。$x=1$ なら $y=-\dfrac{8}{1}=-8$。ほかも\R{同}{おな}じように\R{計算}{けい|さん}する。\par
\noindent\begin{minipage}[c]{0.56\linewidth}
\begin{bktable}{colspec={Q[c,m,wd=6mm] *{8}{Q[c,m,wd=4.6mm]}}, rowsep=1pt, colsep=2pt, cells={font=\footnotesize}}
$x$ & $-8$ & $-4$ & $-2$ & $-1$ & 1 & 2 & 4 & 8 \\
$y$ & 1 & 2 & 4 & 8 & $-8$ & $-4$ & $-2$ & $-1$ \\
\end{bktable}
\par\vspace{4pt}
\kotae \R{表}{ひょう}は\R{左}{ひだり}のとおり。グラフは\R{右}{みぎ}の\R{図}{ず}。$a=-8$ が\R{負}{ふ}なので、\R{左上}{ひだり|うえ}と\R{右下}{みぎ|した}の\R{区画}{く|かく}を\R{通}{とお}る\R{双曲線}{そう|きょく|せん}。
\end{minipage}\hfill
\begin{minipage}[c]{0.4\linewidth}\centering
\begin{tikzpicture}
\begin{axis}[bkplot, width=46mm, height=44mm, xmin=-9.5, xmax=9.5, ymin=-9.5, ymax=9.5, xtick={-8,-4,4,8}, ytick={-8,-4,4,8}, xlabel={$x$}, ylabel={$y$}]
  \addplot[bkcurve, domain=0.86:9.3, samples=100] {-8/x};
  \addplot[bkcurve, domain=-9.3:-0.86, samples=100] {-8/x};
  \addplot[only marks, mark=*, mark size=1.4pt, bkInk] coordinates {(-8,1) (-4,2) (-2,4) (-1,8) (1,-8) (2,-4) (4,-2) (8,-1)};
\end{axis}
\end{tikzpicture}
\end{minipage}
\end{reidai}

\begin{jibun}
$y=\dfrac{6}{x}$ について、$x=-6$、$-3$、$-2$、$-1$、1、2、3、6 のときの $y$ の\R{表}{ひょう}を\R{作}{つく}り、グラフをノートにかきましょう。グラフはどの\R{区画}{く|かく}を\R{通}{とお}りますか。
\bkjans{$y$ は\R{順}{じゅん}に $-1$、$-2$、$-3$、$-6$、6、3、2、1。\R{右上}{みぎ|うえ}と\R{左下}{ひだり|した}の\R{区画}{く|かく}を\R{通}{とお}る\R{双曲線}{そう|きょく|せん}（\R{軸}{じく}にはくっつかない）}
\end{jibun}

\begin{reidai}[グラフから\R{式}{しき}を\R{読}{よ}み、\R{点}{てん}がのっているか\R{調}{しら}べる]
\noindent\begin{minipage}[c]{0.6\linewidth}
\R{右}{みぎ}のグラフは\R{反比例}{はん|ぴ|れい}のグラフで、\R{点}{てん} $(2,\ 5)$ を\R{通}{とお}ります。（1）$y$ を $x$ の\R{式}{しき}で\R{表}{あらわ}しましょう。（2）\R{点}{てん} $(4,\ 2.5)$ と\R{点}{てん} $(-5,\ 2)$ は、このグラフの\R{上}{うえ}にありますか。
\end{minipage}\hfill
\begin{minipage}[c]{0.36\linewidth}\centering
\begin{tikzpicture}
\begin{axis}[bkplot, width=42mm, height=40mm, xmin=0, xmax=10.5, ymin=0, ymax=10.5, xtick={2,4,6,8,10}, ytick={2,4,...,10}, xlabel={$x$}, ylabel={$y$}]
  \addplot[bkcurve, domain=0.97:10.3, samples=100] {10/x};
  \addplot[only marks, mark=*, mark size=1.6pt, bkRed] coordinates {(2,5)};
\end{axis}
\end{tikzpicture}
\end{minipage}\par
\kangaekata グラフの\R{上}{うえ}の\R{点}{てん}は、どれも $xy=a$ をみたします。（2）は、$x\times y$ を\R{計算}{けい|さん}して $a$ と\R{同}{おな}じになるかを\R{調}{しら}べます。\par
\shiki（1）\nobreak $a=2\times5=10$\quad だから $y=\dfrac{10}{x}$\par
\hspace*{2.6\zw}（2）\nobreak $4\times2.5=10$\quad 10と\R{同}{おな}じなので、のっている\par
\hspace*{2.6\zw}\phantom{（2）}$(-5)\times2=-10$\quad 10とちがうので、のっていない\par
\kotae（1）\nobreak $y=\dfrac{10}{x}$\quad（2）\nobreak $(4,\ 2.5)$ はのっている。$(-5,\ 2)$ はのっていない
\end{reidai}

\begin{jibun}
\R{反比例}{はん|ぴ|れい}のグラフが\R{点}{てん} $(3,\ -4)$ を\R{通}{とお}ります。（1）$y$ を $x$ の\R{式}{しき}で\R{表}{あらわ}しましょう。（2）\R{点}{てん} $(-6,\ 2)$ と\R{点}{てん} $(2,\ 6)$ は、このグラフの\R{上}{うえ}にありますか。
\bkjans{（1）\nobreak $a=3\times(-4)=-12$ で $y=-\dfrac{12}{x}$。（2）\nobreak $(-6)\times2=-12$ なので $(-6,\ 2)$ はのっている。$2\times6=12$ なので $(2,\ 6)$ はのっていない}
\end{jibun}

\begin{reidai}[\R{反比例}{はん|ぴ|れい}の\R{変化}{へん|か}の\R{割合}{わり|あい}]
$y=\dfrac{12}{x}$ で、$x$ が1から3まで\R{変}{か}わるときと、3から6まで\R{変}{か}わるときの、\R{変化}{へん|か}の\R{割合}{わり|あい}をそれぞれ\R{求}{もと}めましょう。\par
\kangaekata \R{変化}{へん|か}の\R{割合}{わり|あい}＝$y$ の\R{変化量}{へん|か|りょう}÷$x$ の\R{変化量}{へん|か|りょう}（3-5）。\R{変化量}{へん|か|りょう}は「あと－まえ」で\R{計算}{けい|さん}します。はじめに、\R{区間}{く|かん}の\R{両}{りょう}はしの $y$ を\R{求}{もと}めます。\par
\shiki $x=1$ で $y=12$、$x=3$ で $y=4$、$x=6$ で $y=2$\par
\hspace*{2.6\zw}1→3：$\dfrac{4-12}{3-1}=\dfrac{-8}{2}=-4$\qquad 3→6：$\dfrac{2-4}{6-3}=\dfrac{-2}{3}=-\dfrac{2}{3}$\par
\kotae 1→3は $-4$、3→6は $-\dfrac{2}{3}$（\R{区間}{く|かん}によってちがう）\par
\vspace{2pt}
\noindent\begin{minipage}[c]{0.54\linewidth}
\textbf{グラフで\R{見}{み}る}\par
\R{変化}{へん|か}の\R{割合}{わり|あい}は、グラフの2\R{点}{てん}を\R{結}{むす}んだ\R{直線}{ちょく|せん}の\R{傾}{かたむ}きです（3-5）。1→3の\R{線}{せん}（\R{赤}{あか}）は\R{急}{きゅう}で、3→6の\R{線}{せん}（オレンジ）はゆるやかです。\R{双曲線}{そう|きょく|せん}が、\R{右}{みぎ}へ\R{行}{い}くほどゆるやかになるからです。
\end{minipage}\hfill
\begin{minipage}[c]{0.42\linewidth}\centering
\begin{tikzpicture}
\begin{axis}[bkplot, width=50mm, height=46mm, xmin=0, xmax=7.5, ymin=0, ymax=13.5, xtick={1,3,6}, ytick={2,4,12}, xlabel={$x$}, ylabel={$y$}]
  \addplot[bkcurve, domain=0.9:7.3, samples=100] {12/x};
  \addplot[bkRed, line width=1.2pt] coordinates {(1,12) (3,4)};
  \addplot[bkOrange, line width=1.2pt] coordinates {(3,4) (6,2)};
  \addplot[only marks, mark=*, mark size=1.6pt, bkInk] coordinates {(1,12) (3,4) (6,2)};
  \node[font=\sffamily\scriptsize, text=bkRed!80!black, anchor=west] at (axis cs:2.2,8.6) {$-4$};
  \node[font=\sffamily\scriptsize, text=bkOrange!80!black, anchor=south west] at (axis cs:4.6,3.2) {$-\frac{2}{3}$};
\end{axis}
\end{tikzpicture}
\end{minipage}
\end{reidai}

\begin{jibun}
$y=\dfrac{24}{x}$ で、$x$ が2から4まで\R{変}{か}わるときと、4から8まで\R{変}{か}わるときの、\R{変化}{へん|か}の\R{割合}{わり|あい}をそれぞれ\R{求}{もと}めましょう。
\bkjans{$x=2$、4、8 で $y=12$、6、3。2→4は $\dfrac{6-12}{4-2}=-3$、4→8は $\dfrac{3-6}{8-4}=-\dfrac{3}{4}$}
\end{jibun}

\begin{reidai}[ボイルの\R{法則}{ほう|そく}：\R{注射器}{ちゅう|しゃ|き}の\R{空気}{くう|き}]
\R{注射器}{ちゅう|しゃ|き}に\qty{30}{mL}の\R{空気}{くう|き}を\R{入}{い}れて\R{先}{さき}をふさぐと、\R{圧力}{あつ|りょく}は\qty{100}{kPa}でした。\R{温度}{おん|ど}は\R{変}{か}わらないとします。（1）\R{体積}{たい|せき}を\qty{20}{mL}にすると、\R{圧力}{あつ|りょく}は\R{何}{なん}kPaですか。（2）\R{圧力}{あつ|りょく}を\qty{300}{kPa}にするには、\R{体積}{たい|せき}を\R{何}{なん}mLにすればよいですか。（3）\R{体積}{たい|せき}を $\dfrac{1}{3}$ にすると、\R{圧力}{あつ|りょく}は\R{何倍}{なん|ばい}ですか。\par
\kangaekata ボイルの\R{法則}{ほう|そく}で $pV$ が\R{一定}{いっ|てい}です。はじめの\R{値}{あたい}から $pV$ を\R{求}{もと}め、$p=\dfrac{pV}{V}$、$V=\dfrac{pV}{p}$ で\R{計算}{けい|さん}します。\par
\shiki $pV=100\times30=3000$\par
\hspace*{2.6\zw}（1）\nobreak $p=\dfrac{3000}{20}=150$\qquad（2）\nobreak $V=\dfrac{3000}{300}=10$\par
\hspace*{2.6\zw}（3）\nobreak $\dfrac{1}{3}$ の\qty{10}{mL}では（2）より\qty{300}{kPa}。$300\div100=3$\par
\kotae（1）\nobreak\qty{150}{kPa}\quad（2）\nobreak\qty{10}{mL}\quad（3）\nobreak3\R{倍}{ばい}
\end{reidai}

\begin{jibun}
\R{自転車}{じ|てん|しゃ}の\R{空気入}{くう|き|い}れのつつに、\qty{100}{kPa}の\R{空気}{くう|き}が\qty{40}{mL}\R{入}{はい}っています。\R{先}{さき}をふさいでおしこみます（\R{温度}{おん|ど}は\R{変}{か}わらないとします）。（1）\qty{25}{mL}にすると\R{圧力}{あつ|りょく}は\R{何}{なん}kPaですか。（2）\qty{400}{kPa}にするには\R{何}{なん}mLにすればよいですか。（3）\R{体積}{たい|せき}を $\dfrac{1}{4}$ にすると、\R{圧力}{あつ|りょく}は\R{何倍}{なん|ばい}ですか。（4）\R{圧力}{あつ|りょく}をはじめの5\R{倍}{ばい}の\qty{500}{kPa}にするには、\R{体積}{たい|せき}をはじめの\R{何分}{なん|ぶん}の1にすればよいですか。
\bkjans{$pV=100\times40=4000$。（1）\nobreak $4000\div25=160$ で\qty{160}{kPa}。（2）\nobreak $4000\div400=10$ で\qty{10}{mL}。（3）\nobreak 4\R{倍}{ばい}。（4）\nobreak $4000\div500=8$ で\qty{8}{mL}、$\dfrac{1}{5}$}
\end{jibun}

\begin{reidai}[\R{表}{ひょう}から\R{反比例}{はん|ぴ|れい}を\R{見}{み}つける：てこ]
てこの\R{左}{ひだり}がわにおもりをつるし、\R{右}{みぎ}がわをばねばかりで\R{下}{した}に\R{引}{ひ}いて、つり\R{合}{あ}わせます。\R{支点}{し|てん}からばねばかりまでの\R{距離}{きょ|り} $x$〔cm〕と、ばねばかりの\R{力}{ちから} $y$〔N〕は\R{下}{した}の\R{表}{ひょう}のようになりました。$y$ を $x$ の\R{式}{しき}で\R{表}{あらわ}し、$x=25$ のときの $y$ を\R{求}{もと}めましょう。\par
\noindent\begin{minipage}[c]{0.55\linewidth}\centering
\begin{bktable}{colspec={Q[l,m,wd=13mm] *{4}{Q[c,m,wd=7mm]}}, rowsep=1pt, colsep=3pt, cells={font=\small}}
$x$（cm） & 10 & 20 & 30 & 40 \\
$y$（N） & 6 & 3 & 2 & 1.5 \\
\end{bktable}
\end{minipage}\hfill
\begin{minipage}[c]{0.42\linewidth}\centering
\begin{tikzpicture}[x=0.8mm, y=0.8mm]
  \fill[bkGray!70] (-2.5,-9) -- (2.5,-9) -- (0,-1) -- cycle;
  \fill[bkInk!75, rounded corners=0.6mm] (-22,-1) rectangle (24,1);
  \draw[bkInk!60, line width=0.5pt] (-14,-1) -- (-14,-4);
  \filldraw[fill=bkGray!45, draw=bkGray!90!black, line width=0.6pt, rounded corners=0.8mm] (-17.5,-10) rectangle (-10.5,-4);
  \draw[bkInk!60, line width=0.5pt] (18,-1) -- (18,-4);
  \filldraw[fill=bkBlue!12, draw=bkBlue!70!black, line width=0.6pt, rounded corners=1mm] (15.5,-18) rectangle (20.5,-4);
  \draw[bk arrow, bkRed, line width=1pt] (18,-19) -- (18,-25);
  \draw[{Stealth[length=3pt]}-{Stealth[length=3pt]}, bkBlue, line width=0.6pt] (0,3) -- (18,3);
  \node[font=\sffamily\scriptsize, text=bkBlue!80!black, anchor=south] at (9,3.6) {$x$〔cm〕};
  \node[font=\sffamily\scriptsize, text=bkRed!80!black, anchor=west] at (19.5,-23) {$y$〔N〕};
\end{tikzpicture}
\end{minipage}\par\vspace{3pt}
\kangaekata 「\R{手順}{て|じゅん}」の①〜③のとおりにします。まず、たての\R{組}{くみ}ごとに $x\times y$ を\R{計算}{けい|さん}します。\par
\shiki $10\times6=60$、$20\times3=60$、$30\times2=60$、$40\times1.5=60$\par
\hspace*{2.6\zw}どれも60なので\R{反比例}{はん|ぴ|れい}で、$a=60$\par
\hspace*{2.6\zw}$y=\dfrac{60}{x}$ に $x=25$ を\R{代入}{だい|にゅう}して $y=\dfrac{60}{25}=2.4$\par
\kotae $y=\dfrac{60}{x}$、$x=25$ のとき\qty{2.4}{N}
\end{reidai}

\begin{jibun}
\qty{6}{V}の\R{電源}{でん|げん}にいろいろな\R{抵抗}{てい|こう}をつなぎ、\R{抵抗}{てい|こう} $x$〔Ω〕と\R{電流}{でん|りゅう} $y$〔A〕をはかりました。$y$ を $x$ の\R{式}{しき}で\R{表}{あらわ}し、\qty{12}{\ohm}の\R{抵抗}{てい|こう}をつないだときの\R{電流}{でん|りゅう}を\R{求}{もと}めましょう。
\begin{bkcenter}
\begin{bktable}{colspec={Q[l,m,wd=22mm] *{4}{Q[c,m,wd=11mm]}}, rowsep=1pt, cells={font=\small}}
$x$（Ω） & 2 & 3 & 4 & 6 \\
$y$（A） & 3 & 2 & 1.5 & 1 \\
\end{bktable}
\end{bkcenter}
\bkjans{$x\times y$ はどれも6なので\R{反比例}{はん|ぴ|れい}で $y=\dfrac{6}{x}$（オームの\R{法則}{ほう|そく} $V=RI$ と\R{同}{おな}じ）。$y=\dfrac{6}{12}=0.5$ で\qty{0.5}{A}}
\end{jibun}

%% ---------------- まちがえやすいところ ----------------
\bksection{tsumazuki}{まちがえやすいところをチェック}\label{bk:3-7:tsuma}

\begin{tsumazuki}{グラフを\R{軸}{じく}にくっつける}
\bkng{$y=\dfrac{12}{x}$ のグラフを、\R{右}{みぎ}のはしで $x$ \R{軸}{じく}につなげてかく}
\bkok{$x$ \R{軸}{じく}に\R{近}{ちか}づけるが、くっつけない}
くっついた\R{所}{ところ}では $y=0$ ですが、$x\times0=0$ は12になりません。グラフの\R{上}{うえ}の\R{点}{てん}は、どれも $xy=12$ です。
\end{tsumazuki}

\begin{tsumazuki}{$x=0$ のとき $y=0$ としてしまう}
\bkng{$y=\dfrac{12}{x}$ は、$x=0$ のとき $y=0$ なので、\R{原点}{げん|てん}を\R{通}{とお}る}
\bkok{$x=0$ のときの $y$ はない。\R{原点}{げん|てん}は\R{通}{とお}らない}
$12\div0$ は\R{計算}{けい|さん}できません（1-14）。$y=ax$（\R{比例}{ひ|れい}）なら $x=0$ で $y=0$ ですが、\R{反比例}{はん|ぴ|れい}はちがいます。
\end{tsumazuki}

\begin{tsumazuki}{\R{反比例}{はん|ぴ|れい}なら、いつも\R{右下}{みぎ|さ}がり？}
\bkng{$y=-\dfrac{6}{x}$ も、$x$ がふえると $y$ がへる}
\bkok{$x>0$ では、$x=1$ で $-6$、$x=2$ で $-3$。$y$ はふえる}
$a<0$ のとき、$x>0$ の\R{枝}{えだ}は\R{右上}{みぎ|あ}がりです。\R{表}{ひょう}を\R{作}{つく}って\R{確}{たし}かめよう。
\end{tsumazuki}

\begin{tsumazuki}{\R{変化}{へん|か}の\R{割合}{わり|あい}を $a$ とする}
\bkng{$y=\dfrac{12}{x}$ の\R{変化}{へん|か}の\R{割合}{わり|あい}は、いつも12}
\bkok{\R{区間}{く|かん}ごとにちがう。1→2なら $-6$、2→3なら $-2$}
「\R{変化}{へん|か}の\R{割合}{わり|あい}はいつも $a$」は、1\R{次関数}{じ|かん|すう} $y=ax+b$ のきまりです（3-5）。\R{反比例}{はん|ぴ|れい}では、\R{毎回}{まい|かい}\R{計算}{けい|さん}します。
\end{tsumazuki}

%% ---------------- 物理・工学 ----------------
\begin{butsuri}
\par\addvspace{3pt}\textbf{エンジンのシリンダー}\par\noindent
\begin{minipage}[c]{0.7\linewidth}
\R{自動車}{じ|どう|しゃ}のエンジンでは、シリンダー（つつ）の\R{中}{なか}でピストンが\R{上}{うえ}に\R{動}{うご}いて、\R{空気}{くう|き}とガソリンをおしちぢめてから\R{火}{ひ}をつけます。ふつうのガソリンエンジンでは、\R{体積}{たい|せき}を\R{約}{やく} $\dfrac{1}{10}$ にします。たとえば、\qty{100}{kPa}・\qty{500}{mL}の\R{空気}{くう|き}を\qty{50}{mL}にすると、\R{温度}{おん|ど}が\R{変}{か}わらなければ $pV=100\times500=50000$ なので、
\bkthrhline{$p=\dfrac{50000}{50}=1000$}
で、\qty{1000}{kPa}（10\R{倍}{ばい}）になります。じっさいは、\R{速}{はや}くおしちぢめると\R{空気}{くう|き}の\R{温度}{おん|ど}が\R{上}{あ}がるので、\R{圧力}{あつ|りょく}はもっと\R{大}{おお}きくなります。その\R{計算}{けい|さん}には、3-11で $pV^{1.4}$ が\R{一定}{いっ|てい}という\R{式}{しき}を\R{使}{つか}います。
\end{minipage}\hfill
\begin{minipage}[c]{0.26\linewidth}\centering
\begin{tikzpicture}[x=1mm, y=1mm]
  \fill[bkBlue!12] (0,22) rectangle (16,34);
  \fill[bkGray!40] (-2.5,6) rectangle (0,36) (16,6) rectangle (18.5,36);
  \fill[bkGray!40] (-2.5,36) rectangle (18.5,39);
  \fill[bkInk!65, rounded corners=0.5mm] (0.3,16) rectangle (15.7,22);
  \draw[bkGray!90!black, line width=1.6pt] (8,17) -- (8,2);
  \fill[bkGray!80] (8,2) circle (1.6);
  \draw[bk arrow, bkRed, line width=1pt] (21.5,16) -- (21.5,27);
  \node[font=\sffamily\scriptsize, text=bkBlue!70!black] at (8,28) {\R{空気}{くう|き}};
  \node[font=\sffamily\scriptsize, text=bkInk, anchor=north] at (8,0) {ピストン};
\end{tikzpicture}
\end{minipage}
\end{butsuri}

\begin{butsuri}
\par\addvspace{3pt}\textbf{\R{同}{おな}じ\R{道}{みち}のりの、\R{速}{はや}さと\R{時間}{じ|かん}}\par\noindent
\begin{minipage}[c]{0.6\linewidth}
\R{東京}{とう|きょう}から\R{新大阪}{しん|おお|さか}までの\R{新幹線}{しん|かん|せん}の\R{線路}{せん|ろ}は、\R{約}{やく}\qty{550}{km}です。\R{平均}{へい|きん}の\R{速}{はや}さ $v$〔km/h〕とかかる\R{時間}{じ|かん} $t$ \R{時間}{じ|かん}は、$vt=550$ で\R{反比例}{はん|ぴ|れい}します。\R{平均}{へい|きん}\qty{220}{km/h}なら $550\div220=2.5$ で2.5\R{時間}{じ|かん}、\qty{275}{km/h}なら2\R{時間}{じ|かん}です。\R{時間}{じ|かん}を30\R{分}{ぷん}\R{縮}{ちぢ}めるには、\R{速}{はや}さを\qty{55}{km/h}も\R{上}{あ}げなければなりません。\R{速}{はや}くするほど、\R{同}{おな}じ30\R{分}{ぷん}を\R{縮}{ちぢ}めるのはむずかしくなります。\R{双曲線}{そう|きょく|せん}が\R{右}{みぎ}へ\R{行}{い}くほどゆるやかになるからです。
\end{minipage}\hfill
\begin{minipage}[c]{0.36\linewidth}\centering
\begin{tikzpicture}
\begin{axis}[bkplot, width=48mm, height=44mm, xmin=0, xmax=360, ymin=0, ymax=4.6, xtick={110,220,330}, ytick={1,2,3,4}, xlabel={$v$}, ylabel={$t$}]
  \addplot[bkcurve, domain=125:350, samples=100] {550/x};
  \addplot[only marks, mark=*, mark size=1.6pt, bkRed] coordinates {(220,2.5) (275,2)};
  \draw[bkRed, line width=0.6pt, dash pattern=on 2pt off 1.5pt] (axis cs:220,0) -- (axis cs:220,2.5) (axis cs:275,0) -- (axis cs:275,2);
\end{axis}
\end{tikzpicture}
\end{minipage}

\par\addvspace{3pt}\textbf{\R{電気}{でん|き}の\R{抵抗}{てい|こう}と\R{電流}{でん|りゅう}}\par\noindent
\R{電圧}{でん|あつ}が\R{同}{おな}じなら、オームの\R{法則}{ほう|そく} $V=RI$（2-15）から $RI$ が\R{一定}{いっ|てい}で、\R{電流}{でん|りゅう}は\R{抵抗}{てい|こう}に\R{反比例}{はん|ぴ|れい}します。\R{家}{いえ}のコンセント（\qty{100}{V}）につないだドライヤーの\R{電熱線}{でん|ねつ|せん}は、\R{抵抗}{てい|こう}を\R{小}{ちい}さくするほど\R{大}{おお}きな\R{電流}{でん|りゅう}が\R{流}{なが}れ、たくさん\R{熱}{ねつ}が\R{出}{で}ます。
\end{butsuri}



\begin{motto}[3]
\textbf{はなれるほど、もっと\R{弱}{よわ}くなる：$\dfrac{1}{x^2}$}\par
\noindent\begin{minipage}[c]{0.6\linewidth}
\R{電球}{でん|きゅう}から\qty{1}{m}はなれた\R{所}{ところ}の\R{明}{あか}るさを1とします。\qty{2}{m}はなれると、\R{光}{ひかり}はたても\R{横}{よこ}も2\R{倍}{ばい}に\R{広}{ひろ}がり、\R{同}{おな}じ\R{光}{ひかり}が $2\times2=4$ \R{倍}{ばい}の\R{面積}{めん|せき}に\R{当}{あ}たります。だから\R{明}{あか}るさは $\dfrac{1}{4}$ です。\qty{3}{m}なら $3\times3=9$ \R{倍}{ばい}の\R{面積}{めん|せき}で $\dfrac{1}{9}$ です。\R{明}{あか}るさ $y$ は、\R{距離}{きょ|り} $x$ の2\R{乗}{じょう}に\R{反比例}{はん|ぴ|れい}し、$y=\dfrac{1}{x^2}$ と\R{書}{か}けます。
\end{minipage}\hfill
\begin{minipage}[c]{0.36\linewidth}\centering
\begin{tikzpicture}[x=1mm, y=1mm]
  \fill[bkYellow!70] (0,0) circle (1.6);
  \foreach \d/\s in {10/3, 20/6, 30/9} {\draw[bkGray, line width=0.6pt] (\d,-\s) rectangle (\d+0.01,\s);}
  \fill[bkYellow!25] (0,0) -- (30,9) -- (30,-9) -- cycle;
  \draw[bkYellow!70!black, line width=0.5pt] (0,0) -- (30,9) (0,0) -- (30,-9);
  \foreach \d/\s in {10/3, 20/6, 30/9} {\draw[bkInk, line width=1pt] (\d,-\s) -- (\d,\s);}
  \foreach \d/\t in {10/1, 20/2, 30/3} {\node[font=\sffamily\scriptsize, anchor=north] at (\d,-10) {\qty{\t}{m}};}
  \node[font=\sffamily\scriptsize, anchor=south] at (10,3.6) {1};
  \node[font=\sffamily\scriptsize, anchor=south] at (20,6.6) {4};
  \node[font=\sffamily\scriptsize, anchor=south] at (30,9.6) {9};
\end{tikzpicture}
\end{minipage}\par\vspace{2pt}
\R{電気}{でん|き}を\R{帯}{お}びたもの\R{同士}{どう|し}がおし\R{合}{あ}う・\R{引}{ひ}き\R{合}{あ}う\R{力}{ちから}や、\R{地球}{ち|きゅう}が\R{月}{つき}を\R{引}{ひ}く\R{力}{ちから}も、\R{距離}{きょ|り}の2\R{乗}{じょう}に\R{反比例}{はん|ぴ|れい}します。\R{距離}{きょ|り}が2\R{倍}{ばい}で $\dfrac{1}{4}$、10\R{倍}{ばい}で $\dfrac{1}{100}$ と、$\dfrac{1}{x}$ よりずっと\R{速}{はや}く\R{弱}{よわ}くなります。
\end{motto}

\begin{butsuri}
\par\addvspace{3pt}\textbf{くぎぬき（バール）とてこ}\par\noindent
\begin{minipage}[c]{0.55\linewidth}
くぎぬきは、てこを\R{使}{つか}った\R{道具}{どう|ぐ}です。くぎを\R{抜}{ぬ}くのに\qty{600}{N}の\R{力}{ちから}がいるとします。くぎから\R{支点}{し|てん}までが\qty{3}{cm}なら、てこのつり\R{合}{あ}い（\R{力}{ちから}×\R{支点}{し|てん}からの\R{距離}{きょ|り}が\R{等}{ひと}しい、2-15）から、\R{手}{て}の\R{力}{ちから} $y$〔N〕と、\R{支点}{し|てん}から\R{手}{て}までの\R{距離}{きょ|り} $x$〔cm〕は、
\bkthrhline{$xy=600\times3=1800$}
で、\R{反比例}{はん|ぴ|れい}します。\qty{30}{cm}の\R{所}{ところ}を\R{持}{も}てば\qty{60}{N}（\R{約}{やく}\qty{6}{kg}の\R{物}{もの}を\R{持}{も}つくらい）、\qty{60}{cm}なら\qty{30}{N}ですみます。\R{長}{なが}い\R{柄}{え}の\R{道具}{どう|ぐ}が\R{楽}{らく}なのは、このためです。
\end{minipage}\hfill
\begin{minipage}[c]{0.42\linewidth}\centering
\begin{tikzpicture}[x=0.82mm, y=0.82mm]
  \fill[bkStem!40] (-8,-5) rectangle (46,0);
  % くぎ（床から出た軸と頭）  % noruby
  \draw[bkGray!90!black, line width=1.2pt] (0,-3) -- (0,6);
  \draw[bkGray!90!black, line width=2pt] (-1.8,6.3) -- (1.8,6.3);
  \node[font=\sffamily\scriptsize, anchor=east] at (-2.2,2.5) {くぎ};
  % バール：爪がくぎの頭の下に入り、曲がった所が床に当たる  % noruby
  \draw[bkInk, line width=2.2pt, line cap=round, line join=round] (-2.6,4.6) -- (1.6,4.6) .. controls (3.6,4.6) and (4.2,1.2) .. (5.6,1.1) -- (42,12.5);
  \fill[bkBlue] (0,4.6) circle (0.9);
  \fill[bkRed] (5.4,0.9) circle (0.9);
  \fill[bkGreen!70!black] (42,12.5) circle (0.9);
  \draw[bkBlue!80!black, line width=0.4pt] (-0.6,5.4) -- (-4,11.5);
  \node[font=\sffamily\scriptsize, text=bkBlue!80!black, anchor=south] at (-4.5,11.5) {\R{作用点}{さ|よう|てん}};
  \draw[bk arrow, bkRed, line width=1pt] (42,22) -- (42,14.6);
  \node[font=\sffamily\scriptsize, text=bkGreen!50!black, anchor=south, align=center] at (40,22.5) {\R{力点}{りき|てん}\\\R{手}{て}の\R{力}{ちから} $y$〔N〕};
  \draw[bkRed!80!black, line width=0.4pt] (5.4,-0.5) -- (5.4,-6.5);
  \node[font=\sffamily\scriptsize, text=bkRed!80!black, anchor=north] at (5.4,-6.5) {\R{支点}{し|てん}};
  \draw[{Stealth[length=3pt]}-{Stealth[length=3pt]}, bkBlue, line width=0.6pt] (5.4,-14) -- (42,-14);
  \draw[bkBlue, line width=0.4pt] (42,-15) -- (42,11);
  \node[font=\sffamily\scriptsize, text=bkBlue!80!black, anchor=north] at (24,-14.7) {$x$〔cm〕};
\end{tikzpicture}
\begin{tikzpicture}
\begin{axis}[bkplot, width=46mm, height=34mm, xmin=0, xmax=66, ymin=0, ymax=210, xtick={20,40,60}, ytick={60,120,180}, xlabel={$x$}, ylabel={$y$}]
  \addplot[bkcurve, domain=8.8:64, samples=100] {1800/x};
  \addplot[only marks, mark=*, mark size=1.5pt, bkInk] coordinates {(30,60) (60,30)};
\end{axis}
\end{tikzpicture}
\end{minipage}
\end{butsuri}

\begin{column}{\R{山}{やま}の\R{上}{うえ}で、おかしの\R{袋}{ふくろ}がふくらむわけ}
\noindent\begin{minipage}[c]{0.74\linewidth}
\R{富士山}{ふ|じ|さん}の\R{頂上}{ちょう|じょう}では、まわりの\R{空気}{くう|き}の\R{圧力}{あつ|りょく}は\R{約}{やく}\qty{63}{kPa}で、\R{地上}{ち|じょう}（\R{約}{やく}\qty{101}{kPa}）の6\R{割}{わり}ほどです。\R{閉}{と}じたおかしの\R{袋}{ふくろ}の\R{空気}{くう|き}は、まわりと\R{同}{おな}じ\R{圧力}{あつ|りょく}になるまでふくらみます。\R{地上}{ち|じょう}で\qty{100}{mL}の\R{空気}{くう|き}なら、$pV=101\times100=10100$ なので、\R{頂上}{ちょう|じょう}では $10100\div63$ で\R{約}{やく}\qty{160}{mL}になります（\R{気温}{き|おん}が\R{変}{か}わる\R{分}{ぶん}、\R{少}{すこ}しずれます）。
\end{minipage}\hfill
\begin{minipage}[c]{0.22\linewidth}\centering
\begin{tikzpicture}[x=0.86mm, y=0.72mm]
  % 山（うすい色）  % noruby
  \fill[bkBlue!10] (-1,0) -- (14,52) -- (29,0) -- cycle;
  % 山の上：ぱんぱんの袋、短い矢印が少し  % noruby
  \filldraw[fill=bkOrange!30, draw=bkOrange!80!black, line width=0.6pt, rounded corners=3.5mm] (8,32) rectangle (20,42);
  \fill[bkOrange!70!black] (7,32.8) rectangle (8.2,41.2) (19.8,32.8) rectangle (21,41.2);
  \foreach \a/\b in {{(14,46.5)}/{(14,43.5)}, {(14,27.5)}/{(14,30.5)}, {(2.8,37)}/{(5.8,37)}, {(25.2,37)}/{(22.2,37)}} {\draw[-{Stealth[length=2.5pt]}, bkNavy, line width=0.7pt] \a -- \b;}
  \node[font=\sffamily\footnotesize, anchor=south, align=center, inner sep=1pt] at (14,47) {\R{山}{やま}の\R{上}{うえ}\\\R{約}{やく}\qty{63}{kPa}};
  % ふもと：平らな袋、長い矢印がたくさん  % noruby
  \filldraw[fill=bkOrange!30, draw=bkOrange!80!black, line width=0.6pt, rounded corners=1mm] (8,8.5) rectangle (20,12);
  \fill[bkOrange!70!black] (7,8.8) rectangle (8.2,11.7) (19.8,8.8) rectangle (21,11.7);
  \foreach \xx in {10,14,18} {\draw[-{Stealth[length=2.5pt]}, bkNavy, line width=0.9pt] (\xx,18.5) -- (\xx,13); \draw[-{Stealth[length=2.5pt]}, bkNavy, line width=0.9pt] (\xx,2) -- (\xx,7.5);}
  \draw[-{Stealth[length=2.5pt]}, bkNavy, line width=0.9pt] (0,10.25) -- (5.8,10.25);
  \draw[-{Stealth[length=2.5pt]}, bkNavy, line width=0.9pt] (28,10.25) -- (22.2,10.25);
  \node[font=\sffamily\footnotesize, anchor=north, align=center, inner sep=1pt] at (14,1) {ふもと\\\R{約}{やく}\qty{101}{kPa}};
\end{tikzpicture}
\end{minipage}
\end{column}

%% ---------------- 章末問題 ----------------
\bkendproblems
\bkproblemgroup{\R{基本}{き|ほん}}{1}

\bkproblem $y$ は $x$ に\R{反比例}{はん|ぴ|れい}し、$x=2$ のとき $y=9$ です。$y$ を $x$ の\R{式}{しき}で\R{表}{あらわ}し、$x=-3$ のときの $y$ と、$y=4$ になる $x$ を\R{求}{もと}めましょう（\bkreidairef{1}）。
\bkpans{$a=2\times9=18$ で $y=\dfrac{18}{x}$。$x=-3$ のとき $y=-6$。$4x=18$ で $x=4.5$}

\bkproblem $y=-\dfrac{12}{x}$ について、$x=-6$、$-4$、$-3$、$-2$、2、3、4、6 のときの $y$ の\R{表}{ひょう}を\R{作}{つく}り、グラフをノートにかきましょう。グラフはどの\R{区画}{く|かく}を\R{通}{とお}りますか（\bkreidairef{2}）。
\bkpans{$y$ は\R{順}{じゅん}に2、3、4、6、$-6$、$-4$、$-3$、$-2$。\R{左上}{ひだり|うえ}と\R{右下}{みぎ|した}の\R{区画}{く|かく}を\R{通}{とお}る\R{双曲線}{そう|きょく|せん}}

\bkproblem \R{次}{つぎ}の\R{点}{てん}のうち、$y=\dfrac{20}{x}$ のグラフの\R{上}{うえ}にあるものをすべて\R{選}{えら}びましょう（\bkreidairef{3}）。\nopagebreak\par\nopagebreak
\hspace*{1\zw}ア\enspace $(4,\ 5)$\qquad イ\enspace $(-2,\ -10)$\qquad ウ\enspace $(5,\ -4)$\qquad エ\enspace $(0.5,\ 40)$
\bkpans{ア、イ、エ（$x\times y$ が20になるもの。ウは $-20$）}

\bkproblem $y=\dfrac{18}{x}$ で、$x$ が1から3まで\R{変}{か}わるときと、3から6まで\R{変}{か}わるときの、\R{変化}{へん|か}の\R{割合}{わり|あい}を\R{求}{もと}めましょう（\bkreidairef{4}）。
\bkpans{1→3は $\dfrac{6-18}{3-1}=-6$、3→6は $\dfrac{3-6}{6-3}=-1$}

\bkproblem \R{注射器}{ちゅう|しゃ|き}に\qty{100}{kPa}の\R{空気}{くう|き}を\qty{50}{mL}\R{入}{い}れ、\R{先}{さき}をふさぎました（\R{温度}{おん|ど}は\R{変}{か}わらないとします）。（1）\qty{20}{mL}にすると\R{圧力}{あつ|りょく}は\R{何}{なん}kPaですか。（2）\R{圧力}{あつ|りょく}を\qty{125}{kPa}にするには、\R{何}{なん}mLにすればよいですか（\bkreidairef{5}）。
\bkpans{$pV=100\times50=5000$。（1）\nobreak $5000\div20=250$ で\qty{250}{kPa}。（2）\nobreak $5000\div125=40$ で\qty{40}{mL}}

\bkproblem かみ\R{合}{あ}う\R{歯車}{は|ぐるま}で、\R{歯}{は}の\R{数}{かず} $x$ が10、20、30、40のとき、1\R{分間}{ぷん|かん}の\R{回転数}{かい|てん|すう} $y$ は60、30、20、15でした。$y$ を $x$ の\R{式}{しき}で\R{表}{あらわ}し、\R{歯}{は}が24の\R{歯車}{は|ぐるま}の\R{回転数}{かい|てん|すう}を\R{求}{もと}めましょう（\bkreidairef{6}）。
\bkpans{$x\times y$ はどれも600。だから\R{式}{しき}は $y=\dfrac{600}{x}$。$600\div24=25$ で25\R{回転}{かい|てん}}

\bkproblemgroup{チャレンジ}{2}

\bkproblem \R{反比例}{はん|ぴ|れい}のグラフが\R{点}{てん} $(3,\ 4)$ を\R{通}{とお}ります。このグラフは、\R{点}{てん} $(6,\ b)$ と\R{点}{てん} $(c,\ -2)$ も\R{通}{とお}ります。$b$ と $c$ を\R{求}{もと}めましょう（\bkreidairef{1}・\bkreidairef{3}）。
\bkpans{$a=3\times4=12$ で $y=\dfrac{12}{x}$。$b=\dfrac{12}{6}=2$。$-2c=12$ で $c=-6$}

\bkproblem エンジンのシリンダーで、\qty{100}{kPa}・\qty{500}{mL}の\R{空気}{くう|き}をおしちぢめます（\R{温度}{おん|ど}は\R{変}{か}わらないとします）。（1）\qty{50}{mL}にすると\R{圧力}{あつ|りょく}は\R{何}{なん}kPaですか。（2）\R{圧力}{あつ|りょく}を\qty{1200}{kPa}にするには、\R{体積}{たい|せき}をはじめの\R{何分}{なん|ぶん}の1にすればよいですか。
\bkpans{$pV=100\times500=50000$。（1）\nobreak $50000\div50=1000$ で\qty{1000}{kPa}。（2）\nobreak $1200\div100=12$ \R{倍}{ばい}なので、\R{体積}{たい|せき}は $\dfrac{1}{12}$}

\bkproblem \qty{120}{km}はなれた\R{町}{まち}へ、\R{車}{くるま}で\R{行}{い}きます。\R{平均}{へい|きん}の\R{速}{はや}さを $x$〔km/h〕、かかる\R{時間}{じ|かん}を $y$ \R{時間}{じ|かん}とします。（1）$y$ を $x$ の\R{式}{しき}で\R{表}{あらわ}しましょう。（2）$x=40$、$x=60$ のときの $y$ を\R{求}{もと}めましょう。（3）1.5\R{時間}{じ|かん}で\R{着}{つ}くには、\R{平均}{へい|きん}\R{何}{なん}km/hで\R{走}{はし}ればよいですか。
\bkpans{（1）\nobreak $xy=120$ で $y=\dfrac{120}{x}$。（2）\nobreak $x=40$ で3\R{時間}{じ|かん}、$x=60$ で2\R{時間}{じ|かん}。（3）\nobreak $1.5x=120$ で $x=80$、\qty{80}{km/h}}

\bkproblem \R{地上}{ち|じょう}（\qty{100}{kPa}）で\qty{120}{mL}の\R{空気}{くう|き}が\R{入}{はい}った\R{袋}{ふくろ}を、\R{圧力}{あつ|りょく}が\qty{80}{kPa}の\R{高}{たか}い\R{山}{やま}に\R{持}{も}っていきました。\R{温度}{おん|ど}は\R{変}{か}わらないとすると、\R{袋}{ふくろ}の\R{空気}{くう|き}は\R{何}{なん}mLになりますか（\bkreidairef{5}）。
\bkpans{$pV=100\times120=12000$。$12000\div80=150$ で\qty{150}{mL}}

\bkproblem \qty{12}{V}の\R{電源}{でん|げん}に\R{抵抗}{てい|こう}をつなぐと、\R{抵抗}{てい|こう} $x$〔Ω〕が3、4、6のとき、\R{電流}{でん|りゅう} $y$〔A〕は4、3、2でした。$y$ を $x$ の\R{式}{しき}で\R{表}{あらわ}し、\qty{8}{\ohm}のときの\R{電流}{でん|りゅう}を\R{求}{もと}めましょう（\bkreidairef{6}）。
\bkpans{$x\times y$ はどれも12なので $y=\dfrac{12}{x}$。$y=\dfrac{12}{8}=1.5$ で\qty{1.5}{A}}

\bkproblem $y=-\dfrac{18}{x}$ のグラフが、\R{点}{てん} $(3,\ b)$ と\R{点}{てん} $(c,\ 2)$ を\R{通}{とお}ります。$b$ と $c$ を\R{求}{もと}めましょう（\bkreidairef{1}・\bkreidairef{3}）。
\bkpans{$b=-\dfrac{18}{3}=-6$。$xy=-18$ に $y=2$ を\R{入}{い}れて $2c=-18$、$c=-9$}

\bkproblem \R{先}{さき}をふさいだ\R{注射器}{ちゅう|しゃ|き}の\R{空気}{くう|き}の\R{体積}{たい|せき}を $\dfrac{1}{5}$ にすると、\R{圧力}{あつ|りょく}は\R{何倍}{なん|ばい}になりますか。\R{温度}{おん|ど}は\R{変}{か}わらないとします（\bkreidairef{5}）。
\bkpans{$pV$ が\R{一定}{いっ|てい}なので、\R{体積}{たい|せき}が $\dfrac{1}{5}$ なら\R{圧力}{あつ|りょく}は5\R{倍}{ばい}}

\begin{furikaeri}
\begin{checklist}
\item \R{双曲線}{そう|きょく|せん}をかき、\R{軸}{じく}にくっつかないわけを\R{言}{い}える。\bkfurref{bk:3-7:jiku}
\item $a<0$ のグラフを\R{読}{よ}み、\R{変化}{へん|か}の\R{割合}{わり|あい}を\R{計算}{けい|さん}できる。\bkfurref{bk:3-7:wariai}
\item $xy$ が\R{一定}{いっ|てい}を\R{見}{み}つけ、\R{式}{しき}を\R{作}{つく}って\R{使}{つか}える。\bkfurref{bk:3-7:pv}
\end{checklist}
\end{furikaeri}

\bkfukidashi{sora}{おしちぢめた\R{空気}{くう|き}も、\R{新幹線}{しん|かん|せん}の\R{時間}{じ|かん}も、\R{同}{おな}じ\R{双曲線}{そう|きょく|せん}になるんだね。\R{次}{つぎ}の\R{章}{しょう}では、かっこをはずす\R{計算}{けい|さん}で、ボールのいちばん\R{高}{たか}い\R{点}{てん}を\R{見}{み}つけるよ。}
